\begin{textsample}{POS Dim 3 – 2018\_c6 – Score 1.00 – 2018\_c6/enade\_2018\_tec\_des\_graf\_b78dccf50c.txt}  \label{ex:f3_pos_120}
CAPÍTULO 6 ANÁLISE TÉCNICA DA PROVA Este capítulo tem por objetivo apresentar o desempenho dos estudantes concluintes de Tecnologia em Design Gráfico no Enade/2018.
Para isso, foram calculadas as estatísticas básicas da prova em seu todo ( seção 6.1.1 ), bem como as estatísticas relacionadas aos Componentes de Formação Geral ( seção 6.1.2 ) e de Conhecimento Específico da Área ( seção 6.1.3 ).
Dadas as suas características, foram analisadas, em separado, as questões objetivas ( seção 6.2 ) e as questões discursivas ( seção 6.3 ).
Para as questões objetivas, foram disponibilizados os índices de facilidade e de discriminação Ponto-Bisserial, também em separado, para os Componentes de Formação Geral ( seção 6.2.1 ) e de Conhecimento Específico ( seção 6.2.2 ).
De cada componente, uma das questões foi escolhida para exemplificar a análise gráfica, relacionando -se as alternativas escolhidas pelos estudantes ( inclusive o gabarito ) com o número de acertos no componente.
No Anexo I, apresenta -se a íntegra da análise gráfica para todas as questões objetivas.
Para cada uma das questões discursivas, os conteúdos dos tipos mais comuns de respostas dos estudantes são apresentados e comparados com o padrão de respostas esperado ( ver Anexo VIII com o padrão de respostas ).
Tomando -se como base as duas questões discursivas do Componente de Formação Geral, na seção 6.3.1.6 são apresentados comentários sobre a correção das respostas em relação à Língua Portuguesa.
Nas tabelas, constam as seguintes estatísticas das notas²⁰ : média do desempenho na prova, erro padrão da média, desvio padrão, nota mínima, mediana e nota máxima para cada um de seus componentes.
Tais estatísticas contemplam o total de estudantes concluintes da Área de Tecnologia em Design Gráfico inscritos e presentes à prova do Enade/2018, tendo em vista agregações, ou por Grandes Regiões e o país como um todo, ou por Categoria Administrativa, Organização Acadêmica da IES e Modalidade de Ensino.
Em relação aos gráficos de distribuição de notas, o intervalo considerado foi de 10 unidades, aberto à esquerda e fechado à direita, com exceção do primeiro intervalo, [ 0 ; 10 ], fechado em ambos os extremos.
Para os gráficos de distribuição das notas das questões discursivas, foram consideradas mais duas categorias : questão em branco²¹ e nota zero. ²⁰ Uma definição dessas estatísticas pode ser encontrada no Glossário. ²¹ Nesse grupo estão incluídas também as respostas classificadas como nulas ou desconsideradas. 6.1 ESTATÍSTICAS BÁSICAS DA PROVA Esta seção apresenta estatísticas selecionadas e histogramas da nota geral ( 6.1.1 ) e de cada componente : Formação Geral ( 6.1.2 ) e Conhecimento Específico ( 6.1.3 ).
São também apresentadas estatísticas selecionadas de subpopulações, caracterizadas por Grande Região, Categoria Administrativa, Organização Acadêmica e Modalidade de Ensino. 6.1.1 Estatísticas Básicas Gerais Na Tabela 6.1, são apresentadas as Estatísticas Básicas da prova, por Grande Região, dos estudantes concluintes de Tecnologia em Design Gráfico.
A Média das notas da prova, como um todo ( nas seções seguintes serão analisados os Componentes de Formação Geral e de Conhecimento Específico ), foi 50,0, sendo que os alunos da região Norte obtiveram a Média mais baixa ( 47,1 ), e os da região Sul obtiveram a Média mais alta ( 54,9 ).
As demais médias foram : 49,7 na região Nordeste, 49,3 na região Sudeste e 49,8 na região Centro-Oeste.
O Desvio padrão para o Brasil, como um todo, foi 14,0, sendo o maior Desvio padrão encontrado na região Centro-Oeste ( 14,5 ), e o menor, na região Norte ( 13,2 ), indicando uma dispersão um pouco menor das notas dessa última região.
A região que obteve a maior nota Máxima foi a Nordeste ( 85,3 ), ao passo que a região que atingiu a menor nota Máxima foi a Centro-Oeste ( 80,8 ).
A Mediana do Brasil, como um todo, foi 50,8, sendo a maior Mediana obtida na região Sul ( 55,5 ), e a menor, obtida na região Norte ( 47,5 ).
A nota Mínima do Brasil foi obtida na região Sul, onde foi igual a zero.
A nota Mínima na região Norte foi 16,1, a maior nota Mínima entre as cinco Grandes Regiões.
Considerando -se as notas segundo Grande Região, observa -se que existe diferença estatisticamente significativa ao nível de 95 \% entre a maior Média, obtida na região Sul, e todas as Médias das demais regiões²². ²² Todas as comparações deste capítulo utilizam os intervalos de 95 \%.
Os erros-padrão da média que possibilitam os testes estão disponíveis nas tabelas.
O comportamento das notas dos estudantes de todo o Brasil pode ser observado no Gráfico 6.1, que apresenta um histograma com a distribuição das mesmas : uma distribuição unimodal com a moda no intervalo ( 50 ; 60 ].
Na Tabela 6.2, são apresentadas informações referentes à Média da nota final, desagregadas por Categoria Administrativa, por Organização Acadêmica e por Modalidade de Ensino.
A Média das notas da prova, como um todo, foi 50,0.
Em relação à Categoria Administrativa, os estudantes das IES Públicas obtiveram Média mais alta ( 56,9 ) que a Média nacional, e os das IES Privadas obtiveram Média mais baixa ( 49,2 ).
Observa -se que existe diferença estatisticamente significativa entre as médias das notas das IES Públicas e as das Privadas.
A diferença entre as médias das regiões Sul e Norte ( 7,8 ), a maior e a menor Média, é um pouco superior à diferença entre a das IES Públicas e a das Privadas ( 7,7 ), caracterizando maior diversidade regional do que administrativa.
Os desvios padrão, tanto para as IES Públicas ( 13,4 ), quanto para as IES Privadas ( 13,8 ), foram inferiores ao do Brasil, como um todo ( 14,0 ).
No tocante à Organização Acadêmica, os CEFET/IFET e os Centros Universitários obtiveram Média mais alta que a nacional ( 54,3 e 50,5, respectivamente ).
As médias das Universidades e das Faculdades foram menores do que a nacional ( 49,2 e 49,3, respectivamente ).
Constata -se que existe diferença estatisticamente significativa ao nível de 95 \% entre a Média mais alta, 54,3 dos CEFET/IFET, e a dos demais tipos de Organização Acadêmica.
Cumpre destacar que a Modalidade Educação a Distância não foi representada por nenhum estudante concluínte da Área de Tecnologia em Design Gráfico, portanto, as estatísticas referentes à Modalidade Educação Presencial serão sempre as mesmas que as do Brasil como um todo. 6.1.2 Estatísticas Básicas no Componente de Formação Geral Na Tabela 6.3, são apresentadas as Estatísticas Básicas em relação ao componente da prova que avalia a Formação Geral dos estudantes concluintes.
Os alunos de todo o Brasil obtiveram desempenho médio de 47,6.
Quanto à variabilidade, o Desvio padrão das notas dos estudantes do Brasil, como um todo, foi 18,0.
A maior Média foi obtida na região Sul ( 52,0 ), e a menor, na região Norte ( 42,6 ).
As demais médias foram : 47,3 na região Nordeste, 46,9 na região Sudeste e 48,4 na região Centro-Oeste.
Já o maior Desvio padrão foi obtido na região Sul ( 18,8 ), e o menor, na região Centro-Oeste ( 17,3 ).
Os demais desvios padrões foram : 18,3 na região Norte, 18,1 na região Nordeste e 17,6 na região Sudeste.
A maior nota no Componente de Formação Geral da prova do Enade/2018 foi 95,8, obtida por, pelo menos, um aluno na região Sudeste.
A menor nota Máxima foi obtida na região Norte ( 83,5 ).
A Mediana do Brasil, como um todo, foi 47,5, sendo a menor Mediana encontrada na região Norte ( 40,3 ), e a maior, na região Sul ( 52,5 ).
A nota Mínima nessa parte foi zero em todas as regiões, com exceção da região Centro-Oeste, cuja nota Mínima foi 7,5.
Considerando -se as notas, segundo Grande Região, observa -se que só existe diferença estatisticamente significativa entre a Média das notas do Componente de Formação Geral obtida na região Sul, a maior, e as médias obtidas nas regiões Norte, Nordeste e Sudeste.
As diferenças entre as demais Médias não são estatisticamente significativas.
No Gráfico 6.2, pode -se verificar a avaliação do desempenho dos estudantes do Componente de Formação Geral, a partir do histograma da distribuição das notas correspondentes.
A distribuição é unimodal, com moda em ( 40 ; 50 ], intervalo imediatamente anterior ao modal da distribuição de notas da prova, como um todo ( Gráfico 6.1 ).
Nota -se ainda que, no Gráfico 6.2, as notas apresentam maior dispersão do que as no Gráfico 6.1 ( distribuição das notas da prova ), confirmado pela comparação dos desvios padrões : 14,0 para a nota da prova, como um todo, e 18,0 para o Componente de Formação Geral.
Na Tabela 6.4, são apresentadas as informações referentes ao desempenho dos concluintes do Componente de Formação Geral, em diferentes agregações : Categoria Administrativa, Organização Acadêmica e Modalidade de Ensino.
Observa -se que existe diferença estatisticamente significativa entre as médias dos tipos de Categoria Administrativa : a maior Média obtida por estudantes de IES Públicas ( 56,9 ), e a menor, pelos de IES Privadas ( 46,4 ).
Na desagregação por Organização Acadêmica, a maior Média foi obtida pelos CEFET/IFET ( 52,8 ) e a menor foi a das Faculdades ( 46,1 ).
Nas Universidades a Média foi 46,9 e nos Centros Universitários foi 47,9.
A maior nota Máxima foi obtida em Faculdades ( 95,8 ) e a menor nos CEFET/IFET ( 90,2 ).
Já a maior Mediana foi a dos CEFT/IFET ( 53,3 ), maior que a Mediana nacional ( 47,5 ).
Nota -se que há diferença estatisticamente significativa apenas entre a maior das médias, obtida nos CEFET/IFET ( 52,8 ), e as médias dos demais três tipos de Organização Acadêmica. 6.1.3 - Estatísticas Básicas do Componente de Conhecimento Específico Na Tabela 6.5, são apresentadas as Estatísticas Básicas referentes ao Componente de Conhecimento Específico da área de Tecnologia em Design Gráfico.
A Média do desempenho dos alunos do Brasil, como um todo, foi 50,8.
A maior Média foi obtida na região Sul ( 55,8 ), e a menor, na região Norte ( 48,5 ).
As demais médias foram : 50,4 na região Nordeste, 50,0 na região Sudeste e 50,2 na região Centro-Oeste.
Quanto à variabilidade das notas, o Desvio padrão do Brasil, como um todo, foi 15,0, sendo o maior Desvio padrão observado na região Centro-Oeste ( 15,9 ), e o menor, na região Norte ( 13,8 ).
Os demais desvios foram : 15,0 nas regiões Nordeste e Sudeste, e 13,9 na região Sul.
A Mediana das notas dos estudantes de todo o Brasil foi 51,5.
A maior Mediana ocorreu na região Sul ( 56,8 ), e a menor, na região Norte ( 47,8 ).
As demais medianas foram : 51,3 na região Nordeste, 50,3 na região Sudeste e 51,7 na região Centro-Oeste.
A nota Máxima do Brasil, como um todo, foi 88,0, sendo obtida por, pelo menos, um aluno na região Sudeste.
As demais notas máximas foram : 83,6 nas regiões Norte, 87,3 na região Nordeste, 85,8 na região Sul e 84,1 na região Centro-Oeste.
A nota Mínima foi zero nas regiões Sudeste e Sul, enquanto foi 18,9 na região Norte e 9,4 nas regiões Nordeste e Centro-Oeste.
Observa -se que existe diferença estatisticamente significativa entre a maior Média das notas do Componente de Conhecimento Específico, obtida na região Sul ( 55,8 ), e as médias das demais regiões.
Assim como nos Gráficos 6.1 e 6.2, no Gráfico 6.3, apresentado a seguir, é possível ser feita uma avaliação do desempenho de concluintes de Tecnologia em Design Gráfico em relação ao Componente de Conhecimento Específico, com um histograma da distribuição das notas correspondentes.
Esta também é uma distribuição unimodal, e o grupo modal é o ( 50 ; 60 ], o mesmo o grupo modal da prova, como um todo, e acima do grupo modal para a Formação Geral.
Na Tabela 6.6, é apresentada uma comparação dos resultados em relação à Categoria Administrativa, à Organização Acadêmica e à Modalidade de Ensino, agora levando -se em conta o desempenho de alunos do Componente de Conhecimento Específico da prova.
No que se refere à Organização Acadêmica, a maior Média foi a dos CEFET/IFET ( 54,8 ), vindo a seguir a dos Centros Universitários ( 51,4 ), a das Faculdades ( 50,3 ) e, por fim, a das Universidades ( 50,0 ).
O maior Desvio padrão, e acima do valor para o Brasil, como um todo, foi o das Faculdades ( 15,7 ).
As Faculdades também obtiveram a maior nota Máxima ( 88,0 ).
Os Centros Universitários obtiveram nota Máxima 87,3 ; as Universidades, 86,0 ; e os CEFET/IFET obtiveram nota Máxima 84,8.
As medianas foram : 50,3 nas Universidades, a menor delas ; 52,0 nas Faculdades ; 52,3, nos Centros Universitários ; e 55,9 nos CEFET/IFET, a maior.
A nota Mínima foi zero para todas as Organizações Acadêmicas, com exceção dos CEFET/IFET ( 11,0 ).
Observa -se que existe diferença estatisticamente significativa ao nível de 95 \% no Componente de Conhecimento Específico entre a nota dos CEFET/IFET e as notas dos demais tipos de Organização Acadêmica.
Quanto à Categoria Administrativa, observa -se um comportamento semelhante àquele da parte de Formação Geral e à prova, como um todo, ou seja, existe diferença estatisticamente significativa entre as médias das IES Públicas ( 56,9 ) e as das IES Privadas ( 50,1 ).
Nesse caso também, a maior Média foi obtida por alunos de IES Públicas de ensino. 6.2 ANÁLISE DAS QUESTÕES OBJETIVAS Nesta seção apresentam -se estatísticas selecionadas e histogramas das Questões Objetivas de Formação Geral ( 6.2.1 ) e de Conhecimento Específico ( 6.2.2 ).
São também apresentadas e comparadas as médias das subpopulações, caracterizadas por Grande Região. 6.2.1 Componente de Formação Geral Na Tabela 6.7, são apresentadas as Estatísticas Básicas relativas às oito questões objetivas do componente da prova que abrange a Formação Geral dos estudantes.
A Média do Brasil foi 51,9.
A menor Média foi encontrada na região Norte ( 45,2 ), e a maior, na região Sul ( 57,6 ).
As demais médias foram : 53,7 na região Nordeste, 50,2 na região Sudeste e 55,2 na região Centro-Oeste.
O Desvio padrão do Brasil foi 23,3, sendo o maior Desvio padrão encontrado na região Sul ( 23,4 ), e o menor, na região Nordeste ( 22,7 ).
Os demais desvios foram : 23,2 na região Norte, 23,1 na região Sudeste e 23,3 na região Centro-Oeste.
As medianas do Brasil, como um todo, e de três das cinco regiões foram 50,0.
Na região Norte, a Mediana foi menor, 37,5 ; já na região Sul a Mediana foi maior, 62,5.
A nota Máxima 100,0 e a nota Mínimas zero foram iguais para todas as regiões, sem exceção.
Na Tabela 6.8, são apresentados o Índice de Facilidade e o Índice de Discriminação ( Ponto-Bisserial ) para cada uma das questões objetivas do Componente de Formação Geral.
Quanto ao Índice de Facilidade, foram usadas as seguintes cores para diferenciar o nível de dificuldade da questão : • Azul para as questões classificadas com índice Muito fácil ( \&gt ; =0,86 ), verde para as questões classificadas com índice Fácil ( 0,61 a 0,85 ), amarelo para as questões classificadas com Médio ( 0,41 a 0,60 ), vermelho para as questões classificadas com Difícil ( 0,16 a 0,40 ) e roxo para as questões classificadas com Muito difícil ( \&lt ; =0,15 ).
Já quanto ao Índice de Discriminação, foram usadas as seguintes cores para qualificar a questão : • As questões classificadas com índice Fraco receberam a cor vermelha ( \&lt ; =0,19 ), as classificadas com Médio receberam a cor amarela ( 0,20 a 0,29 ), as classificadas com Bom receberam a cor verde ( 0,30 a 0,39 ), e as classificadas com Muito bom ( \&gt ; =0,40 ) receberam a cor azul.
As questões objetivas do Componente de Formação Geral, segundo o Índice de Facilidade, foram assim avaliadas : das oito questões, nenhuma teve o Índice de Facilidade classificado como Muito fácil, e uma questão foi tida como Fácil, por ter índice de acertos 0,63.
Seis questões foram consideradas com índice de dificuldade Médio, situando -se no intervalo entre 0,41 e 0,60 do Índice de Facilidade, ou seja, houve entre 41,0 \% e 60,0 \% de acertos, enquanto uma questão foi classificada na categoria Difícil, situando -se no intervalo entre 0,16 e 0,40.
Nenhuma questão apresentou menos de 15 \% de acertos, razão pela qual não houve questão classificada como Muito difícil.
O Índice de Facilidade variou de 0,31 a 0,63, e o de Discriminação, de 0,44 a 0,51.
Na Tabela 6.9, é apresentada a distribuição das questões, levando -se em conta simultaneamente a classificação dos dois índices.
As oito questões foram classificadas com Índice de Discriminação Muito bom e, como já comentado, figuraram entre três níveis de dificuldade Fácil, Médio e Difícil : uma classificada na categoria Fácil ( a questão 2 ), seis na categoria Médio ( questões 1, 3, 5, 6, 7 e 8 ), e uma, na categoria Difícil ( questão 4 ).
Em particular, as questões 4 e 5 foram as que apresentaram o maior poder discriminatório, ambas com índice 0,51.
O máximo de acertos foi alcançado pela questão 2, com um Índice de Facilidade de 0,63.
No Gráfico 6.4, para exemplificar, analisa -se o comportamento da questão objetiva de número 4 de Formação Geral.
Trata -se de uma questão considerada Difícil, em relação à facilidade, com índice 0,31, e uma das duas que obteve maior Índice de Discriminação dessa parte da prova ( 0,51 ).
Neste gráfico, cada uma das cinco curvas representa o percentual de respostas em determinada alternativa da questão, em função do número de acertos dos estudantes nessa parte da prova ( Formação Geral/Múltipla Escolha ), antes de possíveis eliminações pelo critério do Ponto-Bisserial.
A curva em verde corresponde à alternativa B, a alternativa correta para essa questão.
Observa -se que, entre os estudantes com menor número de acertos, nessa parte do exame, a situação mais frequente foi a escolha da alternativa C ( em preto ), incorreta.
Por exemplo, entre os estudantes que acertaram três questões, 19,2 \% escolheram a alternativa E ( em vermelho ), 20,6 \% escolheram a alternativa D ( em \textbf{laranja} ), 12,4 \% escolheram a alternativa B ( em verde ), 12,0 \% escolheram a alternativa A ( em azul ), e 35,7 \%, a C ( em preto ).
Entre os que acertaram três respostas nas questões de múltipla escolha de Formação Geral, 0,2 \% deixou a questão em branco e não houve participantes que tenha marcado mais de uma alternativa, invalidando a questão.
À medida que o número de acertos aumenta, indicando desempenho melhor nessa parte da prova, aumenta concomitantemente a proporção de estudantes que selecionaram a alternativa correta B, atingindo 100 \% para os estudantes com oito acertos.
Essa análise permite verificar como a questão discriminou os grupos de desempenho, justificando -se o alto índice obtido na questão.
Cumpre notar que não é possível inferir deste gráfico nem o índice de facilidade ( que seria uma média da proporção ponderada pela quantidade de alunos com cada uma das notas ), nem o índice de discriminação Ponto-Bisserial, por razão equivalente.
No caso extremo, no qual a grande concentração dos acertos dos alunos fosse abaixo de quatro, o índice de facilidade seria obrigatoriamente abaixo de 20 \% ( neste exemplo ).
Caso a concentração fosse em seis acertos ou mais, o índice seria obrigatoriamente acima de 50 \%.
Os gráficos relativos às demais questões de Formação Geral constam do Anexo I. 6.2.2 Componente de Conhecimento Específico Na Tabela 6.10, são apresentadas as Estatísticas Básicas em relação às questões objetivas do Componente de Conhecimento Específico da prova de Tecnologia em Design Gráfico, por Grande Região.
Nessa parte da prova, duas questões foram anuladas pela Comissão Assessora da Área, as questões 19 e 23.
Além disso, como será discutido ainda nesta seção, sete questões objetivas do Componente de Conhecimentos Específicos não foram usadas no cômputo das notas por terem sido descartadas pelo critério do Ponto-Bisserial.
Assim, as notas foram calculadas com base em 18 das 27 questões objetivas de Conhecimentos Específicos.
A Média do Brasil desse componente foi 55,4.
A menor Média foi observada na região Norte ( 53,3 ), e a maior, na região Sul ( 60,5 ).
O Desvio padrão de todo o Brasil foi 16,5, sendo o menor Desvio padrão encontrado na região Sul ( 15,2 ), e o maior, na região Centro-Oeste ( 17,6 ).
A Mediana de todo o Brasil foi 55,6, o mesmo valor da Mediana encontrada em quatro das cinco regiões.
A exceção foi a região Sul ( 61,1 ).
A nota Máxima da prova foi 100,0, obtida nas questões objetivas do Componente de Conhecimento Específico, por, pelo menos, um aluno da região Sudeste.
A nota Mínima foi zero nas regiões Sudeste e Sul.
Na Tabela 6.11, são apresentados os Índices de Facilidade e Discriminação ( Ponto-Bisserial ) das questões objetivas do Componente de Conhecimento Específico, para os estudantes de Tecnologia em Design Gráfico.
Para facilitar a diferenciação das questões, serão usadas as mesmas cores da Tabela 6.8 para as diferentes classificações dos Índices de Facilidade e de Discriminação.
Observa -se que duas questões foram anuladas pela CAA, e que as análises dizem respeito às 25 questões válidas.
A partir do Índice de Facilidade obtido, pode -se concluir que pouco menos da metade das questões objetivas da prova foram considerados, pelo menos, Difícil : das 25 questões válidas, 11 foram classificadas como Difícil ou como Muito difícil.
Uma questão foi classificada como Muito fácil, seis como Fácil e outras sete foram classificadas como Médio.
Já quanto aos índices de discriminação das questões objetivas do Componente de Conhecimento Específico da prova, obteve -se como resultado a seguinte classificação : 12 das 25 questões válidas foram consideradas como boas, enquanto três delas tiveram Índice de Discriminação Muito bom.
Assim, para um pouco menos de dois terços das questões – 15 em 25 – os Índices de Discriminação foram Bom ou Muito bom.
Entre as demais questões, três delas foram classificadas como Médio, e outras sete, como Fraco, sendo dez, por conseguinte, a quantidade de questões nos dois patamares mais baixos de discriminação.
Constata -se, assim, que a prova – no que se refere ao Componente de Conhecimento Específico – possuía média capacidade de discriminar entre aqueles que dominam ou não o conteúdo.
Na Tabela 6.12, é apresentada a distribuição das questões, levando -se em conta simultaneamente a classificação dos dois índices.
Entre as três questões que alcançaram os maiores índices de discriminação, Muito bom, uma delas foi a que teve o maior Índice de Discriminação, a de número 32.
Seu índice foi 0,46, e, quanto ao índice de Facilidade, foi classificada como Fácil, com 66 \% dos estudantes marcando a opção correta.
A questão de número 10 foi a mais difícil entre as 25 questões específicas válidas, com baixo Índice de Facilidade, apenas 14 \% de acertos.
Essa questão apresentou poder discriminatório muito baixo, 0,13, o que comprova ter sido a mais difícil para os estudantes.
Destaca -se, também, a questão 33, com Índice de Facilidade 0,15, o que, em termos percentuais, corresponde a 15 \% de estudantes que responderam acertadamente.
Seu Índice de Discriminação também foi 0,15, Fraco.
Tais questões foram, portanto, pelo critério Ponto-Bisserial, consideradas inadequadas.
Por isso, as questões 10 e 33, além das questões 16, 18, 22, 25 e 28, foram eliminadas do cômputo da nota final.
A título de exemplo das análises do comportamento das questões objetivas, no Gráfico 6.5, analisa -se a questão 32 do Componente de Conhecimento Específico.
Essa questão foi considerada, pelas respostas dos estudantes avaliados na prova, como uma questão com nível de facilidade Fácil, 0,66, ou seja, 66 \% dos estudantes assinalaram acertadamente a opção D, correspondente ao gabarito.
Como já comentado, seu Índice de Discriminação foi igual a 0,46, classificado como Muito bom, o maior valor de discriminação.
Neste gráfico, cada uma das cinco curvas representa o percentual de respostas em determinada alternativa da questão 32, em função do número de acertos dos estudantes nessa parte da prova, antes de possíveis eliminações de questões pelo critério do Ponto-Bisserial.
A alternativa correta D, representada no gráfico pela curva em \textbf{laranja}, foi escolhida em maiores proporções pelos alunos com desempenho melhor nessa parte da prova.
Já as alternativas incorretas, também denominadas distratores, foram selecionadas, principalmente, por aqueles com notas mais baixas.
Observa -se que a soma não é 100 \%, por causa das questões não respondidas ou com mais de uma opção marcada.
Aqueles com nota zero e um, na sua quase totalidade, deixaram essa questão em branco ou marcaram mais de uma alternativa, comportamento considerado inválido.
A proporção de alunos que selecionou a resposta correta D aumenta gradativamente, chegando a atingir 100 \% para 18 acertos, 19 acertos e 21 acertos, enquanto a proporção dos que escolheram alternativas incorretas decai, a partir de sete acertos, em função do número de acertos nessa parte da prova.
Os gráficos relativos às demais questões do Conhecimento Específico constam do Anexo I. 6.3 ANÁLISE DAS QUESTÕES DISCURSIVAS Esta seção apresenta estatísticas selecionadas e histogramas das Questões Discursivas de Formação Geral ( 6.3.1 ) e Conhecimento Específico ( 6.3.2 ).
São também apresentadas e comparadas as médias de subpopulações, caracterizadas por Grande Região. 6.3.1 Componente de Formação Geral As análises dos resultados de desempenho dos estudantes de Tecnologia em Design Gráfico nas duas questões discursivas relativas à Formação Geral encontram -se na Tabela 6.13 e no Gráfico 6.6.
Na Tabela 6.13, observa -se que a nota Média nesse conjunto de questões foi abaixo da obtida nas objetivas.
Os estudantes de todo Brasil obtiveram, em Formação Geral, Média 51,9 nas questões objetivas e 41,1 nas questões discursivas.
Pode -se notar também que o Desvio padrão nesse conjunto de questões foi menor do que o obtido nas objetivas : 23,3 nas questões objetivas e 19,9 nas questões discursivas.
A maior Média foi obtida na região Sul ( 43,6 ), e a menor, na região Nordeste ( 37,7 ).
A Mediana de todo o Brasil, neste componente, foi 43,5.
Nas regiões Norte ( 41,2 ), Nordeste ( 40,5 ) e Centro-Oeste ( 41,8 ), a Mediana foi menor, e nas demais regiões foi maior que o valor para o Brasil : Sudeste ( 44,0 ) e Sul ( 46,5 ).
A nota Máxima ( 90,5 ) foi obtida na região Sudeste, sendo 78,5 a Máxima na região Norte ; 82,0, na Nordeste ; 85,5, na Sul ; e 88,5, na região Centro-Oeste.
A nota Mínima ( 0,0 ) foi a mesma em todas as regiões do Brasil.
No Gráfico 6.6, está representada a distribuição das notas nas questões discursivas do Componente de Formação Geral.
A moda dessa distribuição ocorre no intervalo ( 40 ; 50 ].
Destaca -se também, como um máximo local, o intervalo [ 0 ; 10 ], com distribuição em torno de 10 \% do total de notas, sendo que no intervalo [ 0 ; 10 ] inclui -se, além da nota zero, a frequência de alunos que deixaram esse tipo de questão em branco.
Na sequência, os resultados verificados para cada uma das questões discursivas de Formação Geral serão apresentados, estabelecendo -se relações com os conteúdos abordados em cada uma delas.
Os comentários da Banca de docentes corretores, a respeito do observado na correção das respostas dos estudantes, suas impressões e conclusões serão apresentados juntamente com a análise de cada questão.
Cumpre esclarecer que, tendo -se em vista que as questões discursivas de Formação Geral são padronizadas, ou seja, constam de todas as provas, os comentários da Banca são os mesmos para todas as carreiras acadêmicas, sendo direcionados a todos os estudantes que participaram do Enade/2018.
A seguir, serão analisados os desempenhos da Área de Tecnologia em Design Gráfico nas duas questões discursivas de Formação Geral do Enade/2018, comparando -se os resultados obtidos com comentários para cada questão. 6.3.1.1 Análise de Conteúdo da Questão Discursiva 1 do Componente de Formação Geral Os dados de Tecnologia em Design Gráfico, obtidos a partir das respostas à questão 1, encontram -se na Tabela 6.14 e no Gráfico 6.7.
Nessa questão – a de pior desempenho entre as duas de Formação Geral – os alunos de todo o Brasil tiveram Média 36,4.
A maior Média para a questão 1 foi obtida na região Sul ( 38,2 ), e a menor, na região Centro-Oeste ( 32,9 ).
Quanto à variabilidade das notas, o Desvio padrão de todo o Brasil foi 24,0.
O menor Desvio padrão foi obtido na região Sudeste ( 23,7 ), e o maior foi obtido na região Nordeste ( 25,0 ).
As medianas de três regiões, Norte, Nordeste e Centro-Oeste, foram iguais a 30,0.
Na região Sudeste, a Mediana foi 35,0, a mesma do Brasil como um todo, e na região Sul foi 40,0.
As notas máxima e mínima da questão discursiva 1 foram, sem exceção, as mesmas para todas as regiões do Brasil, respectivamente, 100,0 e zero.
No Gráfico 6.7, mostra -se a distribuição das notas na questão discursiva 1 do Componente de Formação Geral.
Observa -se que a maior frequência corresponde à dos alunos que obtiveram nota situada no intervalo ( 30 ; 40 ], com 17,3 \% dos participantes.
A frequência desse intervalo modal é seguida de perto pelo intervalo ( 40 ; 50 ] e, além disso, o intervalo das questões deixadas em branco se caracteriza como um máximo local, com 11,1 \% dos participantes. 6.3.1.2 Comentários sobre a correção de Conteúdo das respostas à Questão Discursiva 1 A Questão 1 exigia do aluno um posicionamento sobre conteúdo humanista e crítico.
O comando da questão solicitava a redação de um texto que formulasse e articulasse argumentos consistentes envolvendo temas previstos na Portaria nº 444, de 30 de maio de 2018, que dispõe sobre o componente de Formação geral do Enade 2018, tais como : solidariedade/violência, tolerância/intolerância, inclusão/exclusão, sexualidade, relações de gênero e relações étnico-raciais ( art. 7º, item VIII ), bem como ética, democracia e cidadania ( art. 7º, item I ).
O enunciado da questão buscava o estabelecimento da relação de três dimensões – os Direitos Humanos, a Democracia e as Fake News.
Para isso, recorreu a quatro suportes : a reprodução de uma foto e três fragmentos adaptados de textos que foram disponibilizados em páginas da Internet ( da Anistia Internacional, do Jornal El País-Brasil e do Labic-Laboratório de estudos sobre Imagem e Cibercultura ).
Os fragmentos traziam dados sobre o aumento do número de assassinatos de grupos minoritários, as origens e o trabalho da vereadora Marielle Franco, sua luta pelos Direitos Humanos, seu assassinato, e a produção de fake news. vereadora Marielle Franco, sua luta pelos Direitos Humanos, seu assassinato, e a produção de fake news.
O comando solicitava um texto abordando dois tópicos : o tensionamento entre a defesa dos Direitos Humanos e a produção de notícias falsas, tendo como base a vereadora Marielle – a defesa dos Direitos Humanos realizada por ela e as notícias falsas divulgadas após a sua morte.
As fake news também estavam presentes no segundo tópico de análise : os prejuízos que provocam para a sociedade democrática.
Por um lado, os textos selecionados para constituir o enunciado da questão continham apoio suficiente para formulação e articulação de argumentos para discorrer sobre o tensionamento entre a atuação da Marielle em defesa dos Direitos Humanos e fake news.
Por outro, havia ênfase nos aspectos biográficos e na trajetória da vereadora Marielle, o que levou a respostas pouco focadas no comando, que discutiam a vida e a morte da vereadora, sua importância como um ícone dos Direitos Humanos, opiniões em torno do seu assassinato, até mesmo sobre a responsabilidade pelo crime, muitos textos com a presença/defesa de fake news ou produção de novas.
O foco sobre Marielle também ensejou respostas basicamente centradas na questão da violência.
Foram frequentes as interpretações com viés político-partidário, não previstas no padrão de resposta.
O item do comando que solicitava a discussão sobre democracia e fake news encontrou menos apoio nos textos.
Apesar de se esperar que graduandos tivessem maturidade, inclusive política, para demonstrar alguma compreensão dos aspectos principais que envolvem o convívio democrático, esse aspecto foi pouco discutido e, quando presente, quase sempre foi apenas mencionado, numa quase reprodução do comando.
A falta de entendimento do que é o conceito de Direitos Humanos também foi percebida nas respostas, seja para recomendá -lo como importante para a sociedade, seja para criticá -lo como algo que defende bandidos.
Foi marcante a ideia que vincula Direitos Humanos a direitos apenas das minorias.
Não foi incomum encontrar respostas que personificavam os Direitos Humanos, tratando -os como sujeito de alguma ação que deveria ser implementada na sociedade em geral ou em casos específicos.
As fake news foram citadas muitas vezes, mas pouco explicadas, conforme previsto no padrão de resposta.
O argumento utilizado geralmente vinculava as consequências de fake news para a imagem da Marielle, ou seja, centrava no fato em si, já indicado no enunciado, sem qualquer análise ou aprofundamento próprio.
É importante salientar que muitas respostas, provavelmente de grupos vinculados a cursos específicos, continham argumentações recorrentes : os Direitos Humanos como algo vinculado à Constituição e a tratados dos quais o Brasil é signatário ; as fake news como produto do mundo digital e globalizado, relacionado à era da pós-verdade ; a discussão sobre o papel das mídias sociais e da imprensa. produto do mundo digital e globalizado, relacionado à era da pós-verdade ; a discussão sobre o papel das mídias sociais e da imprensa.
Como já apontado, houve baixa quantidade de respostas contemplando o segundo tópico do comando – prejuízos da produção de notícias falsas para a sociedade democrática ; fato motivado, possivelmente, pela falta de um elemento que ressaltasse esse ponto nos textos motivadores.
De modo geral, os estudantes não conseguiram recorrer a conhecimentos tácitos sobre democracia para construir argumentos próprios.
As respostas que abordavam o tema, geralmente estavam vinculadas à ideia do direito à verdade, do direito à integridade da pessoa ou dos direitos gerais garantidos pela Constituição.
A relação da questão com o momento político brasileiro foi indicada em um número expressivo de respostas.
Por vezes, apenas questionando a temática da questão – caso Marielle – em uma prova para avaliar o ensino superior no Brasil ou a importância dos Direitos Humanos na atual conjuntura política do país, levando à anulação da resposta.
Houve também interpretações das consequências do uso das fake news no processo eleitoral de 2018, influenciando a escolha consciente e livre das pessoas.
As respostas demonstraram atenção político social frente a esses temas, porém com pouco domínio conceitual do campo de conhecimento sobre Direitos Humanos, fake news e democracia.
Em termos da linguagem adotada pelos alunos, há dois grandes grupos de respostas.
No primeiro grupo, os que responderam de forma clara, com autonomia de pensamento e explicitação de justificativas para os aspectos solicitados no enunciado, demonstrando conhecimento sobre o que significam Direitos Humanos, fake news e democracia, além da articulação entre essas questões.
Por outro lado, um segundo grupo dos que não conseguiram ir além do que foi citado no enunciado, apenas repetindo -o, muitas vezes desviando o foco da questão para temas correlatos, como a vida de Marielle ou a violência.
Nesse segundo grupo, mais facilmente foram encontradas respostas desconexas e mal articuladas, com textos confusos que explicitam o não conhecimento sobre o tema e a falta de adequação à linguagem, necessária a um concluinte do ensino superior.
As abordagens dos Direitos Humanos mais recorrentes nas respostas corretas foram : a referência da Marielle como representante/ícone dos DH ; e os DH como direitos previstos na Constituição e em tratados internacionais dos quais o Brasil é signatário.
No entanto, também foram encontradas críticas aos Direitos Humanos, tais como : DH vistos como “ direitos de bandidos ” ; a defesa pouco embasada de frases de efeito como “ direitos humanos para humanos direitos ”.
Também chamou atenção o tratamento dos DH como um personagem, ou um órgão público, evidenciando desconhecimento sobre o tema.
Em relação às fake news, dois argumentos previstos pelo padrão de respostas foram os mais presentes : a noção de que o mundo globalizado e a facilidade de reprodução de informações pelas mídias atuais levam à criação de mentiras que são compartilhadas e que geram consequências ; a Marielle como alvo de fake news que tentaram desconstruir sua luta e seu legado.
Também foi frequente a ideia de que a produção de informações falsas é realizada por uma imprensa tendenciosa com intenção de manipular a informação.
No item relacionado a fake news e democracia, as respostas corretas afirmavam que as fake news inviabilizam o acesso à verdade e à garantia da integridade de um Estado Democrático de Direito.
Também foi abordada a influência das fake news nos resultados das eleições dos últimos anos, nos Estados Unidos e no Brasil, gerando consequências para a vida das pessoas.
Por fim, um outro aspecto abordado nas respostas a este tópico foi a polarização da sociedade em termos políticos, com a afirmação de que as fake news dificultam o debate, estimulam posições extremas e rivalidades prejudiciais à convivência em uma democracia plena.
A notas mais fracas ( de zero a 35 ) foram cerca de 20 \% das respostas corrigidas – as que não estavam em branco nem foram desconsideradas ou anuladas.
Foram respostas que se restringiam a cópia ou paráfrases do enunciado ; as que tinham foco em desdobramentos políticos e policiais do caso Marielle ; as que tinham argumentação confusa e/ou incompleta.
A grande maioria das notas, em torno de 65 \%, foi mediana, entre 40 a 70.
Nesse grupo ficaram as respostas que continham alguma abordagem adequada sobre Direitos Humanos, sobre as atividades da vereadora Marielle e sobre as fake news divulgadas após sua morte.
As notas mais altas deste grupo foram para respostas que incluíram a questão do tensionamento, quase sempre, em um texto coerente e consistente.
Muitas vezes, a nota mediana foi consequência do baixo índice de respostas à segunda parte do comando, que solicitava a discussão sobre os prejuízos das fake news para a democracia.
As melhores notas só foram obtidas por aqueles que incluíram a discussão sobre os prejuízos das fake news para a democracia.
Apenas 15 \% das respostas corrigidas receberam avaliação igual ou superior a 75.
A deficiência principal detectada pelos corretores foi a dificuldade de grande parte dos alunos para interpretar o enunciado da questão e para se expressar em textos consistentes.
De modo geral, os temas foram tratados de forma fragmentada e pouco analítica.
Além disso, chamou à atenção a predominância do lugar comum na construção argumentativa das respostas, a superficialidade do debate em torno da democracia e o desconhecimento de conceitos como Direitos Humanos, esses confundidos com instituições e até pessoas. 6.3.1.3 Análise de Conteúdo da Questão Discursiva 2 do Componente de Formação Geral A Tabela 6.15 mostra que o desempenho médio dos estudantes na questão discursiva 2 ( média 39,0 ) foi superior ao obtido na questão discursiva 1 ( média 36,4 ).
A região Sul foi aquela cuja Média, também nessa questão, foi maior ( 41,7 ), enquanto a de menor Média foi a região Norte ( 35,5 ).
Quanto à variabilidade das notas, o Desvio padrão de todo o Brasil foi 25,4, superior ao obtido na questão discursiva 1 ( 24,0 ).
O menor desvio nessa questão foi obtido na região Sudeste ( 24,5 ), enquanto o maior foi obtido na região Sul ( 27,3 ).
A Mediana de todo o Brasil foi 35,0, a mesma de três regiões.
As exceções foram as regiões Sudeste e Sul, que obtiveram medianas 40,0 e 45,0, respectivamente.
A nota Máxima 100,0 foi obtida por, pelo menos, um estudante nas regiões Sudeste e Sul.
A nota Mínima ( 0,0 ) foi a mesma em todas as regiões, sem exceção.
No Gráfico 6.8, mostra -se a distribuição das notas na questão discursiva 2 do Componente de Formação Geral.
Observa -se que a maior frequência corresponde à dos alunos que obtiveram nota situada no intervalo ( 30 ; 40 ].
Destacam -se, também, os intervalos ( 50 ; 60 ] e os das respostas deixadas em branco, com frequências 16,5 \% e 10,7 \% respectivamente.
Nota -se, ainda, que as notas ficaram mais dispersas em comparação às da questão discursiva de número 1, o que pode ser constatado também pela comparação do desvio padrão das notas da questão discursiva 2 ( 25,4 ) e o das notas da questão discursiva 1 ( 24,0 ). 6.3.1.4 Comentários sobre a correção de Conteúdo das respostas à Questão Discursiva 2 Os textos e o comando da Questão 2 de Formação Geral são perfeitamente adequados ao que se espera avaliar em Formação Geral, ao tratar das funções sociais dos museus, contextualizando a temática sob o ponto de vista da memória e da perspectiva de futuro.
Na questão, estavam presentes, principalmente, o conteúdo humanista e crítico, a interpretação e a produção de textos sobre temas como cultura e arte, educação, ciência e democracia, previstos na Portaria que dispõe sobre o componente de Formação Geral do Enade 2018.
Os conhecimentos exigidos para responder ao comando, explicitados no padrão de resposta, vão além do senso comum, por se exigir a compreensão de conceitos como “ função social ” ( de instituições sociais ), “ memória ” ( no sentido de memória social ), “ perspectiva de futuro ” ( no contexto sócio-histórico-cultural ) e “ sociedade contemporânea ” ( que deveria independer da interpretação de cada grupo social a respeito da sociedade na qual vive ).
No entanto, apesar de envolver tais conceitos, espera -se que alunos concluintes de cursos superiores de graduação sejam capazes de discutir funções de instituições sociais, especialmente porque o enunciado trazia fato recente e de grande repercussão como apoio. superiores de graduação sejam capazes de discutir funções de instituições sociais, especialmente porque o enunciado trazia fato recente e de grande repercussão como apoio.
O padrão de resposta elencou as funções sociais que poderiam ser abordadas nas respostas dos estudantes, como também argumentos que poderiam ser usados ao longo das mesmas.
A tarefa dos estudantes era a construção de um texto expositivo que articulasse a importância de museus, suas funções sociais e as perspectivas de memória e de futuro, o que se mostrou de dificuldade de média a difícil.
Apesar de o enunciado ser claro, muitos estudantes, provavelmente influenciados pelos textos, limitaram -se a discutir aspectos relacionados ao Museu Nacional e, em particular, ao incêndio que o destruiu.
Foram poucos os que interpretaram corretamente o escopo do comando, que solicitava a discussão das funções sociais de museus de uma forma geral, no sentido de argumentar a importância deste tipo de instituição.
As quatro funções sociais elencadas no padrão de resposta foram encontradas na quase totalidade das redações dos estudantes de forma articulada com a importância dos museus.
Entretanto, quase nunca os textos abordavam mais do que uma função social e vinculavam, de forma clara, a importância desse tipo de instituição para a sociedade contemporânea.
Além disso, a solicitação de contextualização sob o foco museu/memória/perspectiva de futuro parece não ter sido bem compreendida e, em consequência, esteve pouco presente nas respostas.
Também se observou confusão entre os conceitos de história e memória e uma visão estreita do significado de sociedade contemporânea, às vezes, restrita à cidade do respondente, outras, ao Rio de Janeiro, ou ainda a um tipo de sociedade idealizada.
A relevância que os textos de apoio deram ao caso particular do incêndio do Museu Nacional do Rio de Janeiro e aos prejuízos causados ao seu acervo, induziu a maioria dos estudantes a elaborar respostas tratando das funções a partir da perda e, em particular, à perda do tipo de acervo daquele museu e do tipo de pesquisa que lá se desenvolve, divergindo do comando da questão, que solicitava texto abordando funções sociais de museus em geral.
Também foram encontradas muitas respostas que se limitavam a discutir o que o estudante considerava como possíveis causas do incêndio, a falta de manutenção, questões políticas e de gestão, sem responder ao que foi solicitado.
Foi também recorrente a transcrição de frases dos textos do enunciado nas respostas, diretamente ou por paráfrases, revelando certa dificuldade de criação textual.
A análise das respostas dos estudantes, além de revelar dificuldade de abordar a temática esperada, revelou outros dois aspectos predominantes.
Um deles, como já comentado, foi a ênfase no incêndio do Museu Nacional, que orientou grande parte dos textos dos estudantes.
Outro, também evidenciando desvio do esperado, foi o foco na perda patrimonial.
Nesse grupo, houve respostas que até apresentaram uma função social, mas olhada pelo ponto de vista da perda.
Nessa linha, muitas respostas enalteciam a importância do Museu Nacional, sua relevância no cenário nacional, o impacto da perda de seu acervo, isso complementado por frágeis reflexões sobre as funções sociais.
Em consequência, a perspectiva de desdobramentos mais ricos, que tratassem de aspectos sociais relacionados a diferentes tipos de museu foi prejudicada.
Muitas respostas trataram da preservação da memória e/ou guarda do acervo ; identidade, tanto individual quanto coletiva ; conhecimento sobre antepassados ; evolução da sociedade e dos seus costumes ; valorização e disseminação do conhecimento por meio da liberação do acesso aos acervos ; bem como atividades pedagógicas e exposições de arte.
O aspecto produção de pesquisa e conhecimento, embora frequente, foi mencionado de forma muito simples, sem qualquer aprofundamento, para além do que já era mencionado no enunciado.
Frequentemente apresentava -se o museu como produtor de pesquisas para o campo da saúde, o que evidencia desconhecimento das reais possibilidades de pesquisa, mesmo no Museu Nacional.
Quanto a perspectiva de futuro, quando presente no texto, era mencionada de forma breve, com afirmações genéricas sobre a possibilidade de melhoria do futuro pelo simples conhecimento do passado, numa percepção de que o passado determina o futuro, de forma linear.
Foram raras as respostas com alguma argumentação que associasse pesquisa, memória, identidade e reflexão social como base para a construção de alternativas de futuro.
A relação entre turismo e museu apareceu considerando a visita a museus como forma de lazer e, algumas vezes, como geradora de recursos.
O desempenho pode ser considerado fraco.
Muitas respostas deixavam claro que o estudante nunca tinha visitado um museu e, de uma forma geral, observou -se desconhecimento sobre suas funções sociais.
Não se esperava que história, memória social e museologia fossem assuntos dominados pelos estudantes das diversas áreas que participaram do Enade 2018, mas esperava -se que essa temática pudesse ser abordada do ponto de vista das experiências pessoais ou de conhecimento social.
Como uma vivência própria com os museus também se revelou frágil ou inexistente, o conhecimento sobre o assunto ficou bem restrito aos textos motivadores e ao senso comum, construído, basicamente, pelo noticiário em torno do incêndio.
O universo dos argumentos dos estudantes se mostrou muito limitado.
Cada dimensão foi praticamente representada por uma ideia simplista.
De forma recorrente : a preservação da memória para construção da identidade de um povo ; a compreensão do passado por meio das pesquisas como oportunidade para melhorar o futuro ; o museu como um espaço em que há um aprendizado estimulado pelo lazer.
Houve, também, uma grande concentração de respostas associando o museu a um local de guarda de acervo e de história.
As respostas à Questão 2 evidenciaram dificuldades conceituais.
Além dos conceitos de história, memória e museu se misturarem, parecendo não se distinguirem, ao mencionarem a pesquisa, também se observou falta de distinção entre pesquisas colegiais, com objetivos de aprendizado, e as realizadas por pesquisadores, com objetivos de avançar cientificamente.
A falta de compreensão do que se pedia no comando da questão e a dificuldade de expressão linguística foram dois aspectos marcantes na avaliação dos corretores.
Foi frequente os textos apresentarem ideias soltas ou de forma confusa, sem constituir unidade.
O vocabulário usado se revelou bem restrito, e fortemente apoiado no coloquial, na linguagem oral.
Predominou o senso comum em lugar de uma reflexão mais formal envolvendo os conceitos envolvidos. 6.3.1.5 Análise de Língua Portuguesa das Questões Discursivas do Componente de Formação Geral Os dados de Tecnologia em Design Gráfico, obtidos a partir das respostas às questões discursivas do Componente de Formação Geral, no que tange à Língua Portuguesa, encontram -se na Tabela 6.16 e no Gráfico 6.9.
Nesse aspecto, os alunos de todo o Brasil obtiveram Média 55,0.
A maior Média em relação à Língua Portuguesa foi obtida na região Sul ( 58,5 ), e a menor, na região Nordeste ( 50,0 ).
Quanto à variabilidade das notas, o Desvio padrão de todo o Brasil foi 21,8.
O menor Desvio padrão foi obtido na região Sudeste ( 20,3 ) ; já o maior Desvio padrão foi obtido na região Centro-Oeste ( 24,7 ).
A Mediana das notas de Língua Portuguesa foi 60,0 para o Brasil, como um todo, a mesma obtida nas regiões Norte, Sudeste e Centro-Oeste.
A Mediana da região Sul foi 65,0, a maior entre as obtidas pelas Grandes Regiões, enquanto a da região Nordeste foi 57,5, a menor.
A nota Máxima para todo o Brasil foi 90,0, com, pelo menos, um aluno tirando essa nota nas regiões Sudeste e Sul.
Na região Norte, a nota Máxima foi 82,5 ; na região Nordeste, foi 85,0 ; e na região Centro-Oeste, 87,5.
Já a nota Mínima foi zero em todas as regiões do país.
No Gráfico 6.9, mostra -se a distribuição das notas de Língua Portuguesa do Componente de Formação Geral.
Observa -se que a maior frequência ( quase 30 \% ) corresponde à dos alunos que obtiveram nota no intervalo ( 60 ; 70 ].
Destacam -se, também, os alunos que deixaram as respostas para ambas as questões em branco, representando um pouco mais do que 5 \% do total, caracterizando -se como máximo local. 6.3.1.6 Comentários sobre a correção das respostas de Formação Geral com respeito à Língua Portuguesa As duas questões discursivas de Formação Geral ensejaram duas configurações textuais distintas, tendo em vista os seus encaminhamentos : a Questão 1 teve um caráter mais opinativo e a Questão 2, mais expositivo, conforme se explicita a seguir.
Ambas as questões permitiram a observação do desempenho linguístico dos participantes em situação formal, competência exigida no processo de formação universitária.
As habilidades de produção textual envolvidas têm, igualmente, como condição básica, o domínio do padrão formal culto da Língua Portuguesa, no que diz respeito aos aspectos relativos às questões de caráter ortográfico, textual, morfossintático e vocabular.
O encaminhamento das questões evidencia a expectativa de que o participante utilizasse seus conhecimentos sobre os dois temas e estruturasse seus textos de acordo com as características do registro formal adequado à situação comunicativa.
Tal configuração determina exigências quanto : à adequação da seleção vocabular, ao desenvolvimento do conteúdo, à estruturação sintática dos períodos, à organização lógica das ideias, à utilização de procedimentos de encadeamento textual e de referenciação, à obediência às exigências morfossintáticas próprias da modalidade escrita da norma-padrão, ao respeito às regras ortográficas e às regras de acentuação gráfica.
Esse encaminhamento exigiu do graduando o domínio do modo de organização textual de base dissertativo-argumentativa, essencial ao processo de formação universitária.
As habilidades de produção textual envolvidas têm como condição básica o domínio do padrão formal culto da Língua Portuguesa, no que diz respeito aos aspectos relativos às questões de caráter ortográfico, textual, morfossintático e vocabular.
O enunciado da Questão 1 tomava como base três fragmentos de textos jornalísticos publicados em sites da internet, que abordavam o aumento do assassinato de minorias e de defensores de direitos humanos no Brasil, entre eles a vereadora Marielle Franco.
O comando solicitava que o participante desenvolvesse um texto dissertativo e argumentativo com as temáticas Fake News, Democracia e Direitos Humanos.
Embora não tenha havido referência explícita à tipologia textual esperada para o desenvolvimento da resposta, essa proposta encaminhou o participante para a elaboração de um texto opinativo, com o objetivo de discutir o problema e manifestar um posicionamento crítico.
O enunciado da Questão 2 abordava o incêndio do Museu Nacional do Rio de Janeiro, ocorrido em 2018, e seus impactos na ciência brasileira e internacional, tomando como base dois fragmentos de textos jornalísticos publicados em sites da internet.
O comando solicitava que o participante desenvolvesse um texto abordando três aspectos da função social dos museus na sociedade contemporânea, encaminhando -o para a elaboração de um texto expositivo que confirmasse a posição defendida nos fragmentos motivadores a respeito da importância dessas instituições.
Para atender a essas exigências, o padrão de resposta utilizado na avaliação das questões discursivas de Formação Geral considerou aspectos relevantes ao bom desempenho linguístico como competências distintas, de modo a permitir um mapeamento detalhado do domínio dos recursos disponíveis na Língua Portuguesa para a comunicação escrita formal : ortográficos, textuais, morfossintáticos e vocabulares ( ver detalhes no Anexo VIII – Padrão de Respostas ).
Aspectos Ortográficos Para avaliar a competência relativa ao domínio das convenções ortográficas da norma-padrão da Língua Portuguesa, observou -se o respeito às regras de acentuação gráfica e da grafia padrão das palavras ( com ausência de abreviaturas próprias da linguagem da internet ), de acordo com as convenções estabelecidas pela legislação em vigor e consubstanciadas no Vocabulário Ortográfico da Língua Portuguesa, editado pela Academia Brasileira de Letras.
Aspectos Textuais Para avaliar a competência relativa ao domínio dos procedimentos de estruturação textual, procurou -se observar a utilização de mecanismos linguístico-discursivos que estabeleçam o encadeamento lógico entre as partes do texto, de forma a garantir a progressão e a coerência textuais.
Esta competência destaca -se por sua relevância, já que são avaliadas a clareza e a unidade textual.
Aspectos Morfossintáticos Para avaliar o domínio dos diferentes aspectos morfossintáticos próprios da modalidade escrita formal da norma-padrão da Língua Portuguesa, observou -se se o texto produzido atendeu às seguintes exigências : a concordância nominal, a concordância verbal, a regência nominal, a regência verbal, a flexão nominal, a flexão verbal, a correlação entre os tempos verbais, a colocação pronominal e a utilização de sinais de pontuação que contribuam para a organização lógica da frase e do texto.
Aspectos Vocabulares Para avaliar a adequação da seleção vocabular à modalidade escrita formal da Língua Portuguesa, exigida pela situação comunicativa, observou -se o atendimento às seguintes exigências : precisão na seleção/utilização do vocabulário relacionado à temática solicitada pela questão ; ausência de marcas da oralidade, como termos de sentido muito genérico ( “ coisa ”, “ negócio ”, “ você ” ) e termos de registros mais informais ( como gírias, jargões, frases feitas, ditados populares, termos regionais ). ( “ coisa ”, “ negócio ”, “ você ” ) e termos de registros mais informais ( como gírias, jargões, frases feitas, ditados populares, termos regionais ).
Assim, espera -se que o participante respeite a adequação vocabular não usando gírias ou expressões coloquiais, evite repetição desnecessária de palavras e utilize um vocabulário mais formal, como solicitado por um texto dissertativo.
Para efeito de pontuação, a avaliação do desempenho linguístico considerou três grandes grupos de competências, agrupando aspectos explicitados, com os seguintes pesos relativos : aspectos ortográficos ( 20 \% ) ; aspectos textuais ( 40 \% ), aspectos morfossintáticos e vocabulares ( 40 \% ).
Os critérios de pontuação utilizados no padrão de resposta procuraram atender à diversidade de desempenho dos graduandos.
Nos textos analisados, observaram -se, assim, vários níveis de atendimento às exigências do padrão formal próprio da modalidade escrita da Língua Portuguesa, configurando um"continuum".
Por um lado, textos que traduzem com exatidão a expectativa de uso formal da língua, em todos os níveis analisados, marcados por clareza e unidade ; e, no outro extremo, textos quase incompreensíveis devido ao grande número de desvios de todos os tipos, principalmente os de caráter textual.
As duas questões discursivas de Formação Geral analisadas revelaram desempenhos distintos dos participantes.
A Questão 1 propiciou a oportunidade de maior desenvolvimento do tema solicitado, gerando consequências na estruturação textual, já que os textos foram mais longos, mais elaborados e fluentes.
Em virtude da polêmica sobre a existência das fake news e sobre a visão pública da personagem da vereadora Marielle Franco, observaram -se respostas que polemizaram o tema, favorecendo a construção argumentativa do texto.
A Questão 2 teve um comportamento distinto, devido ao caráter mais expositivo do enunciado e ao reduzido conhecimento dos participantes sobre o tema.
Assim, os textos foram mais curtos, frequentemente com repetição de ideias dos textos motivadores.
Quanto aos aspectos linguísticos analisados durante esta avaliação, observaram -se os seguintes resultados : Aspectos ortográficos : O desempenho dos participantes revelou uma diferença muito grande nos dois aspectos analisados nesta competência : baixo índice de desvios da grafia padrão e grande índice de desvios de acentuação.
Em vários casos, ocorre ausência completa de acentuação gráfica em todas as palavras do texto.
Os resultados revelaram que a tendência dominante entre os universitários brasileiros é a eliminação da acentuação gráfica, talvez motivada pelos hábitos relacionados às redes sociais e pela ausência de esclarecimento dos meios de comunicação, das autoridades e das escolas sobre as decisões do Acordo Ortográfico da Língua Portuguesa de 1990.
Os casos mais sistemáticos de eliminação do acento indicador da sílaba tônica foram : - palavras proparoxítonas ( por exemplo “ generos ”, “ arbitro ”, “ politico ”, “ publicas ”, “ numero ” ) ; - palavras paroxítonas terminadas em ditongo crescente ( por exemplo “ area ”, “ noticiarios ”, “ individuo ”, “ historia ”, “ varios ”, “ noticia ”, “ negligencia ”, “ ocorrencia ” ) ; - palavras oxítonas ( por exemplo “ ninguem ”, “ esta ”, “ ate ”, “ tambem ” ).
Por outro lado, destaca -se o uso indevido do acento gráfico em palavras como “ melâncolia ” e “ intervênção ”, por exemplo.
Quanto ao domínio das convenções relativas à grafia das palavras, observaram -se desvios esporádicos como : hipercorreção pela escolha de “ e ” no lugar de “ i ”, por influência de hábitos da oralidade ( “ descriminação ” no lugar de “ discriminação ”, “ entervenção ” no lugar de “ intervenção ” ) ; desvios de grafia relacionados à variação diastrática, como por exemplo “ perca ”, “ subjulgadas ”, “ precoseito ”, “ soubre ”, “ vecendo ”, “ indesencia ”, “ apolojia ”, “ fulga ”, “ dereitos ” ;...
Observaram -se, também, muitos casos de inadequação no uso da maiúscula : uso indevido para destacar determinadas palavras-chave do texto, como “ Brasileiros ”, “ Homens ” ; ausência de maiúscula para grafar nomes próprios, como"marielle".
Expressões como"Direitos Humanos"foram grafadas diversificadamente, com maiúsculas e com minúsculas.
Destaque -se, também, a maioria dos participantes que grafam os textos inteiramente em caixa alta, o que impede a identificação de desvios relacionados ao uso de maiúsculas.
Vale observar, também, que, ao contrário do que se esperava, não apareceram abreviaturas próprias do “ internetês ”, relacionadas ao uso de redes sociais e e-mails, como por exemplo"vc".
Aspectos textuais : Para alguns dos graduandos essa competência se revelou como a mais problemática, tendo em vista os inúmeros problemas observados, desvios acumulados durante toda a formação escolar.
São eles : sequência justaposta de ideias sem encaixamentos sintáticos, provocando uma redução de estruturas subordinadas, ao lado do aumento na frequência de estruturas coordenadas e absolutas ; frases fragmentadas que comprometem a estrutura lógico-gramatical ; frases formadas apenas por oração subordinada, sem oração principal ; redução no uso de conectores para expressar relações lógicas essenciais à construção do texto como consequência da mudança de estruturação frasal ; emprego inadequado do pronome relativo ( com omissão da preposição ou a utilização do pronome relativo inadequado, como “ onde ” ), refletindo hábitos da oralidade ; ausência de recursos de referenciação, como substituição de termos por sinônimos, hiperônimos, nominalizações, expressões metafóricas.
Esses problemas revelam dificuldades em relação à estrutura formal do texto produzido, o que é preocupante ao se levar em conta que são graduandos em fase final de formação.
Em alguns textos, falta um mínimo de textualidade e de domínio do registro padrão da língua.
Na verdade, observaram -se relações linguísticas quase agramaticais, como as estabelecidas pela sequência de gerúndios sem o apoio de um ponto de partida para a organização das informações gramaticais e semânticas.
Quanto à utilização dos sinais de pontuação, observou -se uma ausência dos mesmos na maioria dos textos analisados.
Foi muito frequente a ocorrência de parágrafos sem marca interna de pontuação para separar os períodos.
São os seguintes os tipos de problemas mais encontrados :  vírgula : utilização de vírgula para separar o sujeito e o predicado ; uso de vírgula no lugar do ponto para separar ideias que constituem períodos distintos ; ausência de vírgula para separar oração adjetiva explicativa ;  ponto e vírgula : utilização do ponto e vírgula no lugar de vírgula ;  ponto final : ausência de ponto final para separar períodos.
Aspectos morfossintáticos e vocabulares : Os resultados são transparentes em relação aos aspectos mais problemáticos no desempenho dos participantes nestas duas competências.
Em relação aos aspectos morfossintáticos, seguem algumas observações sobre os desvios mais frequentes.
Quanto à regência, o desvio mais frequente foi a falta do sinal indicativo da crase – isso revela que o usuário não tem consciência de que, sob a forma do termo “ a ”, existe a presença de uma contração entre a preposição “ a ” ( exigida pela regência do termo anterior ) e o artigo definido “ a ”.
Outro problema relacionado à regência verbal e à nominal, encontrado frequentemente nas respostas, foi a ausência de preposição antes de pronome relativo, processo generalizado lógico-gramatical ; frases formadas apenas por oração subordinada, sem oração principal ; redução no uso de conectores para expressar relações lógicas essenciais à construção do texto como consequência da mudança de estruturação frasal ; emprego inadequado do pronome relativo ( com omissão da preposição ou a utilização do pronome relativo inadequado, como “ onde ” ), refletindo hábitos da oralidade ; ausência de recursos de referenciação, como substituição de termos por sinônimos, hiperônimos, nominalizações, expressões metafóricas.
Esses problemas revelam dificuldades em relação à estrutura formal do texto produzido, o que é preocupante ao se levar em conta que são graduandos em fase final de formação.
Em alguns textos, falta um mínimo de textualidade e de domínio do registro padrão da língua.
Na verdade, observaram -se relações linguísticas quase agramaticais, como as estabelecidas pela sequência de gerúndios sem o apoio de um ponto de partida para a organização das informações gramaticais e semânticas.
Quanto à utilização dos sinais de pontuação, observou -se uma ausência dos mesmos na maioria dos textos analisados.
Foi muito frequente a ocorrência de parágrafos sem marca interna de pontuação para separar os períodos.
São os seguintes os tipos de problemas mais encontrados :  vírgula : utilização de vírgula para separar o sujeito e o predicado ; uso de vírgula no lugar do ponto para separar ideias que constituem períodos distintos ; ausência de vírgula para separar oração adjetiva explicativa ;  ponto e vírgula : utilização do ponto e vírgula no lugar de vírgula ;  ponto final : ausência de ponto final para separar períodos.
Aspectos morfossintáticos e vocabulares : Os resultados são transparentes em relação aos aspectos mais problemáticos no desempenho dos participantes nestas duas competências.
Em relação aos aspectos morfossintáticos, seguem algumas observações sobre os desvios mais frequentes.
Quanto à regência, o desvio mais frequente foi a falta do sinal indicativo da crase – isso revela que o usuário não tem consciência de que, sob a forma do termo “ a ”, existe a presença de uma contração entre a preposição “ a ” ( exigida pela regência do termo anterior ) e o artigo definido “ a ”.
Outro problema relacionado à regência verbal e à nominal, encontrado frequentemente nas respostas, foi a ausência de preposição antes de pronome relativo, processo generalizado na modalidade oral da língua, em situações de registro informal.
Apesar da possibilidade de que essa alteração de regência se generalize no padrão escrito da Língua Portuguesa, como já está ocorrendo até em textos jornalísticos, o não emprego da preposição foi considerado como desvio neste processo de avaliação.
Outro desvio de regência significativo foi a utilização inadequada de uma preposição ou sua ausência após o verbo ou o nome ( substantivo ou adjetivo ).
A concordância verbal e a concordância nominal apresentaram alguns desvios muito frequentes.
Quanto à concordância de número, observou -se, como apontam as pesquisas nessa área, ausência de marca ( com sujeito anteposto ou posposto ) ou uso indevido ( uso inadequado da marca de plural comandado pelo núcleo plural da locução adjetiva, apesar de o substantivo que funciona como núcleo do sintagma nominal estar no singular ).
Uma ocorrência que se destacou foi a ausência de acento circunflexo na forma plural do presente do indicativo do verbo “ ter ”, que foi considerada como um desvio na concordância verbal e não na acentuação gráfica.
Quanto à concordância de gênero, vários casos foram observados, normalmente no âmbito de sintagmas nominais longos, em que o adjetivo está afastado do substantivo.
Quanto à questão da colocação pronominal, foram poucos os desvios observados.
Concluiu -se que, no registro escrito formal, a maioria dos participantes já incorporou regras como a não introdução da frase por um pronome oblíquo e a próclise na presença de um termo atrator.
Várias marcas de oralidade foram identificadas, embora não com alta frequência : uso do pronome relativo “ onde ” como relativo universal, falta de artigo definido antes de substantivo isolado ou antes de termos em enumerações.
Quanto aos aspectos vocabulares, alguns tipos de inadequação foram observados : expressões da oralidade ; seleção vocabular incompatível com o contexto, gerando falta de inteligibilidade ; falta de domínio de vocabulário mais abstrato e de maior complexidade, essencial ao desenvolvimento do texto de base dissertativa.
O principal aspecto observado foi a excessiva repetição de certas palavras, como o termo"pessoa", por exemplo, revelando limitação de repertório vocabular.
Em função do tema solicitado na Questão 1, diferentes grafias da expressão fake news não foram consideradas.
Para sintetizar, os problemas que mais se destacaram nos textos analisados foram : a ) truncamentos sintáticos, caracterizados por estruturas frasais incompreensíveis devido à ruptura da complexidade sintática própria da modalidade escrita ; b ) redução de estruturas subordinadas, compensada pelo aumento na frequência de estruturas coordenadas e absolutas, por um lado, ou pela elaboração de estruturas textuais prejudicadas pelo excesso de ideias sem a devida conexão subordinativa ; c ) redução no uso de conectores para expressar relações lógicas essenciais à construção do texto, substituídos pela exigência de inferência por parte do interlocutor para suprir a sua ausência ; d ) redução no uso do subjuntivo, ao lado da ampliação do uso do indicativo, combinado a estruturas frasais coordenadas ou absolutas ; d ) empobrecimento do processo de referenciação, com a repetição de palavras, sem processos mais sofisticados de substituição ; e ) simplificação da marcação da categoria tempo na morfologia verbal ; f ) desrespeito às regras de concordância verbal e nominal próprias do padrão formal ; g ) ausência do acento grave indicador da crase, revelando falta de conhecimento dos critérios de regência verbal e nominal ; h ) falta de domínio de vocabulário mais abstrato e de maior complexidade, essencial ao desenvolvimento do processo dissertativo ; i ) redução do emprego da acentuação gráfica.
Como resultado da correção, merece atenção a constatação de que a modalidade escrita tem apresentado a tendência a uma evidente simplificação, aproximando -se das características da modalidade oral da Língua Portuguesa.
No caso dos textos de base dissertativa ( expositiva ) e de base argumentativa, inscritos em um registro formal, a distância entre as duas modalidades é ainda maior, o que provoca desvios recorrentes em todos os quatro aspectos analisados : ortográficos, textuais, morfossintáticos e vocabulares.
Destaca -se, nessa avaliação, o comprometimento textual, por meio de estruturas fragmentadas e/ou truncadas, rompendo a complexidade sintática esperada no padrão formal. 6.3.2 Componente de Conhecimento Específico Na parte da prova relativa às questões discursivas do Componente de Conhecimento Específico ( Tabela 6.17 ), observa -se que a Média foi bem mais baixa do que para as questões discursivas do Componente de Formação Geral.
Enquanto no Componente de Formação Geral, a Média para estudantes de Tecnologia em Design Gráfico de todo o Brasil foi 41,1, na parte de Conhecimento Específico, a Média foi 25,1.
A maior Média desse componente foi obtida pelos estudantes da região Sul ( 29,4 ), e a menor, pelos da região Norte ( 21,9 ).
Quanto Quanto à variabilidade das notas, o Desvio padrão de todo o Brasil foi 15,7.
O maior Desvio padrão foi encontrado nas regiões Norte e Sul ( 16,8 ), e o menor, na região Nordeste ( 14,9 ).
A maior nota Máxima, 88,3, foi obtida na região Sudeste.
Nas demais regiões, as notas Máximas foram : 78,3 nas regiões Norte e Nordeste, 85,0 na Sul e 68,3 na região Centro-Oeste.
A nota Mínima ( 0,0 ) foi obtida por, pelo menos, um aluno em todas as regiões do Brasil.
A Mediana do Brasil, como um todo, foi 25,0.
A maior Mediana foi obtida na região Sul ( 30,0 ), e a menor, na região Norte ( 18,3 ).
No Gráfico 6.10, representa -se a distribuição das notas nas questões discursivas no Componente de Conhecimento Específico.
A moda dessa distribuição ocorre no intervalo ( 20 ; 30 ], com 23,1 \% do total de participantes, seguida de perto pelos intervalos [ 0 ; 10 ] e ( 10 ; 20 ].
A análise de cada uma destas questões será feita a seguir. 6.3.2.1 Análise da Questão Discursiva 3 do Componente de Conhecimento Específico Na questão 3, cujos resultados aferidos encontram -se descritos na Tabela 6.18, a Média dos estudantes de todo o Brasil foi 22,1.
A menor Média nessa questão foi obtida pelos alunos das regiões Nordeste e Centro-Oeste ( 19,4 ), enquanto a maior Média foi obtida na região Sul ( 26,1 ).
Quanto à variabilidade das notas, o Desvio padrão de todo o Brasil foi 24,3.
O maior Desvio padrão foi obtido na região Norte ( 27,4 ), enquanto o menor foi obtido na região Nordeste ( 23,0 ).
A nota Máxima, 100,0 pontos, foi alcançada por, pelo menos, um aluno de todas as regiões, com exceção da região Centro-Oeste que teve nota Máxima 85,0.
A Mediana do Brasil, como um todo, foi 10,0, a mesma para as regiões Norte, Nordeste e Sudeste.
Na região Sul, a Mediana foi maior, 20,0, e na região Centro-Oeste foi menor, 0,0.
A nota Mínima ( 0,0 ) também foi a mesma em todas as regiões do Brasil.
Cabe salientar que pelo menos a metade dos alunos na região Centro-Oeste tirou nota zero nessa questão.
No Gráfico 6.11, mostra -se a distribuição das notas na questão discursiva 3, do Componente de Conhecimento Específico da área de Tecnologia em Design Gráfico.
Sem se levar em consideração que 24,3 \% dos estudantes deixaram a resposta em branco e que 16,6 \% receberam nota zero, a moda dessa distribuição seria o intervalo ( 30 ; 40. 6.3.2.2 - Comentários sobre as respostas à Questão Discursiva 3 Na questão discursiva 3, fez -se alusão ao filme “ Pantera Negra ” ( 2018 ), em que foram abordadas as questões raciais, as divisões territoriais e a diversidade cultural dos povos africanos, além de ser retratado o cenário de guerras e de instabilidade política do continente africano.
No enunciado da questão, mencionava -se também que a história e a temática do filme são contemporâneas devido não somente à evolução tecnológica, mas também à reverência à origem e à ancestralidade cultural e religiosa dos negros africanos.
No enunciado também se considerava que o Design se apropria de elementos e narrativas culturais no desenvolvimento de projetos e solicitava -se que o aluno descrevesse uma produção de design gráfico em que os elementos visuais revelassem a apropriação da riqueza étnica e da cultura dos povos africanos.
No enunciado abordam -se pontos importantes das diretrizes curriculares ao se \textbf{traçarem} relações entre a atividade projetual e os aspectos históricos, culturais e sociais para a escolha de referências e elementos visuais.
Para essa escolha, é valorizada a percepção das características e necessidades de grupos sociais, comunidades e indivíduos, correlacionando -as ao ponto mais essencial do fazer profissional do tecnólogo em design gráfico : o ato projetual em si, dando ênfase ao adequado recurso aos elementos visuais originários desses contextos e a devida justificativa para sua escolha como repertório visual em projetos gráficos.
A questão tem o mérito de abordar a cultura africana, conteúdo de caráter obrigatório nacionalmente em todos os cursos, inclusivo e ainda pouco lecionado de forma específica e direta.
O nível de conhecimento exigido para uma boa resposta a essa questão é o esperado de um profissional com formação tecnológica superior, uma vez que para a concepção de projetos de design gráfico sempre é necessário selecionar, articular, justificar e fundamentar a origem histórico-cultural dos elementos visuais que constarão dos trabalhos, o que deve ser utilizado para o convencimento dos clientes de sua adequação e do quanto são vinculados ao vocabulário formal do público ao qual se destina o projeto.
Esse era o ponto focal para que a resposta fosse considerada correta.
No comando da questão, demandou -se claramente a descrição de “ uma produção ” de design gráfico “ em que os elementos visuais revelassem a apropriação da riqueza étnica e da cultura dos povos africanos ”.
Para realizar corretamente a descrição e a exemplificação do referenciamento na “ cultura dos povos africanos ”, o estudante deveria compreender o continente africano como uma fonte de diversidade cultural muito ampla, e como cita o enunciado, partir de uma reflexão sobre as implicações culturais da diáspora africana em outros continentes, em especial nas Américas.
Após essa reflexão, o estudante deveria selecionar um exemplo, dentre os muitos existentes, e elaborar as conexões entre a apropriação de elementos visuais de uma dada cultura de raiz africana e sua correta aplicação ao projeto citado.
Cabe ressaltar que o comando foi bem formulado e claro, e havia referência direta ao que era mais importante : o estudante relacionar para atingir uma boa solução para o problema proposto.
Dentre as três questões discursivas propostas nesta edição do Enade, a Questão 3 pode ser considerada de grau intermediário de dificuldade para sua correta solução, pois apesar de mais focada em aspectos abstratos do exercício profissional, foi elaborada sobre atividade pertinente e basilar a todo e qualquer projeto de design gráfico.
Como dito anteriormente, no texto introdutório da questão utilizou -se de uma narrativa ficcional de uma produção cinematográfica de 2018 ( Pantera Negra ) para ilustrar a temática.
Como consequência, um volume significativo de estudantes estruturou suas respostas a partir da descrição dos elementos visuais do filme, das peças gráficas para sua divulgação, com diferentes graus de aproveitamento.
Houve respostas com base em outros produtos audiovisuais ficcionais ( em especial as versões de Rei Leão ).
Nesses casos, o principal equívoco foi tratar essas produções como se se reportassem literalmente à realidade total do continente africano, sem a percepção de conteúdos estereotipados e generalidades presentes nessas produções.
Embora o comando da questão fosse claro e adequado quanto ao avaliar a capacidade dos estudantes em perceber as ligações necessárias entre as questões culturais e os projetos de design gráfico, parte significativa dos estudantes se limitou a parafrasear o texto do enunciado, enquanto outra parte fez alusão direta ao conteúdo da Questão objetiva 5, que abordou a exposição ExÁfrica.
Em linhas gerais, os estudantes poderiam responder à questão abordando os seguintes itens : aspectos das culturas originárias desse continente que envolvem civilizações da Antiguidade ; a experiência cultural do norte do continente, – fortemente marcada pelo Islamismo – ; ou a profunda diversidade de culturas da África Subsaariana e das populações existentes na bacia do rio Nilo.
Ainda assim, não foram encontradas muitas respostas com diversidade e originalidade.
Um número inexpressivo de respostas abordou itens que não eram relacionados à África subsaariana ou que continham estereótipos de africanidade.
Nem todas apresentaram propostas de projetos ou fizeram conexões com conceitos de História da Arte.
Em parte significativa das respostas, os estudantes apresentaram vocabulário pouco correlacionado ao tema e foram evasivos, não fazendo conexões concretas com projetos gráficos existentes ou hipotéticos.
Contudo, a maioria dos estudantes reduziu as possibilidades de resposta ao conjunto de dados mais presentes na cultura de massa, aos estereótipos da africanidade subsaariana ou a dados das religiões de matriz africana, raramente indo além de citar a umbanda ou o candomblé.
Também foi frequente o uso de vocabulário dúbio em relação ao preconceito racial, como, por exemplo, as expressões"mãe preta", “ mulher de cor ” e “ povo preto ”, de forma desconexa com o campo do design gráfico.
O exemplo mais equivocado que foi encontrado tentava vincular aspectos da cultura contemporânea com a citação do desenho animado Simpsons como exemplo de apropriação dos desdobramentos da cultura dos povos africanos na América.
Nele o estudante se referiu especificamente ao personagem Apu Nahasapeemapetilon, ignorando que o personagem foi declaradamente concebido em aspectos hindus.
Houve uma quantidade significativa de respostas em que os estudantes se reportaram ao campo da moda, mais fortemente ao da estamparia e ao da criação de padrões ; a maioria dando como exemplo influências dos elementos visuais das culturas da África Subsaariana, citando ainda acessórios e bijuterias, e representando através desses elementos o feminino africano, embora muitas vezes de forma estereotipada.
Houve, também, uma quantidade significativa de respostas em que havia referências a projetos ( existentes ou hipotéticos ), como demandado no enunciado.
Nesse campo, a comunicação visual do Mundial de Futebol da FIFA da África do Sul de 2010, o material gráfico do filme Pantera Negra e a exposição ExÁfrica são os exemplos concretos de análise mais utilizados, alguns com um pouco mais de sucesso na abordagem dos aspectos das culturas africanas presentes nos trabalhos, outros com desenvolvimentos distintos da argumentação.
Uma vez que o comando da questão é bastante amplo, o estudante poderia redigir uma resposta correta abordando diferentes aspectos das possibilidades de interpretação do que sejam as culturas africanas, ou suas manifestações no continente americano e, dessa forma, discorrer de muitas formas sobre como elas influenciam a visualidade dos projetos gráficos.
Poderia, por exemplo, citar o uso de elementos da arquitetura e da pintura do Egito Antigo, como as formas geométricas das construções dessa civilização, além dos aspectos da lei da frontalidade e de outras normas da pintura mural egípcia ; trabalhar questões sobre a caligrafia e a influência mulçumana no norte africano, tratar da criação de padrões geométricos, tendo como base os desenhos comuns na arquitetura e no mobiliário dessa região ; escolher aspectos visuais das muitas culturas da África subsaariana, riquíssima em produções visuais bastante veiculadas em espaços públicos, como museus e centros culturais, ou mesmo em matérias de revistas ( digitais e/ou impressas ) especializadas em geografia humana ou cultura.
Justamente nesse aspecto foram identificadas muitas deficiências em relação ao desenvolvimento das argumentações para a construção de respostas adequadas a essa questão.
Chamou a atenção o desconhecimento da diversidade cultural, artística, étnica e geográfica fora do eixo das culturas dominantes na América do Norte e na Europa, além do repertório restrito no que tange à criação de respostas com propostas de projetos hipotéticos.
Foi perceptível o baixo domínio do léxico do campo teórico da área.
Tornou -se evidente para a banca a redução do continente africano a uma versão estereotipada de determinadas regiões do sul da África, com as respostas muitas vezes restritas ao uso de determinados conjuntos de cores, texturas ( baseadas em pelos de animais, em vegetais ou em artesanato ), padrões geométricos e figuras desenhados de forma mais livre.
Também foi muito frequente a referência às estampas estereotipadas como típicas das culturas africanas ou das populações afrodescendentes das Américas.
Outra dificuldade enfrentada foi a compreensão do se exigia e a relação da aplicação dos elementos em um projeto gráfico, concretamente existente ou hipotético, algo pensado pelo estudante no momento da resposta.
O que pode ser considerado grave, uma vez que o curso é para formação de profissionais capazes exatamente de pensar em projetos gráficos novos e diversificados.
Foi constatada uma grande dificuldade com o vocabulário profissional mais cotidiano e com o desenvolvimento de respostas em que articulasse com desenvoltura a associação entre teoria e prática.
Chamou atenção que, em muitos aspectos, as respostas dos estudantes estavam no limiar entre o politicamente incorreto, o desconhecimento cultural e o racismo.
Foi constatado também o significativo número de vezes em que termos equivocados com relação à população afrodescendente foram utilizados na construção do raciocínio dos estudantes.
Outro fato que chama a atenção é a necessidade de um trabalho mais frequente para o esclarecimento da delimitação do campo de atuação do designer gráfico, que fica claro na medida em que muitas respostas não conseguiram se reportar a um projeto que seja tipicamente do campo, além de ter sido muito frequente o recurso aos exemplos originários dos campos do cinema e da moda.
Uma recomendação válida para os cursos da área de Tecnologia em Design Gráfico é que exercícios de interpretação de texto e de redação sejam desenvolvidos em diferentes disciplinas, tendo como temática a apropriação de elementos culturais na concepção dos projetos mais pertinentes ao campo do design gráfico, como identidades visuais, sinalização, criação de tipografias, cartazes, embalagens, dentre outros.
Em se tratando de recomendações para a formação da área em questão, não é possível deixar de relatar a grande dificuldade dos alunos na articulação escrita e na expressão do pensamento de forma textual, com muitas respostas excessivamente curtas, outras ainda em que os alunos apresentaram um grau elevado de dificuldade na conexão dos conteúdos listados, que muitas vezes não possuíam muita conexão nem apresentavam relações claras.
No Gráfico 6.12, está representada a distribuição de notas da questão discursiva 4, do Componente de Conhecimento Específico.
Essa distribuição tem moda no intervalo ( 20 ; 30 ].
Destaca -se o intervalo das respostas deixadas em branco, com 7,3 \% dos estudantes, que se caracteriza como um máximo local. expressão do pensamento de forma textual, com muitas respostas excessivamente curtas, outras ainda em que os alunos apresentaram um grau elevado de dificuldade na conexão dos conteúdos listados, que muitas vezes não possuíam muita conexão nem apresentavam relações claras. 6.3.2.3 - Análise da Questão Discursiva 4 do Componente de Conhecimento Específico Na Tabela 6.19, constam as informações relativas à questão 4 do conjunto do Componente de Conhecimento Específico.
O desempenho dos estudantes de todo o Brasil nesta questão foi o mais alto entre as três questões discursivas desse componente.
A Média geral do Brasil foi 31,1, sendo a maior Média registrada na região Sul ( 35,5 ), e a menor, na região Norte ( 27,2 ).
A nota Máxima ( 100,0 ) foi atingida por, pelo menos, um concluinte das regiões Norte, Nordeste, Sudeste e Sul.
No Centro-Oeste, a nota Máxima foi 70,0.
A Mediana em todo o Brasil foi 30,0, o mesmo valor obtido das regiões Nordeste, Sudeste e Centro-Oeste.
Na região Norte, a Mediana foi menor, 20,0, e na região Sul, a Mediana foi maior, 40,0.
A nota Mínima ( 0,0 ) foi a mesma em todas as regiões.
No Gráfico 6.12, está representada a distribuição de notas da questão discursiva 4, do Componente de Conhecimento Específico.
Essa distribuição tem moda no intervalo ( 20 ; 30 ].
Destaca -se o intervalo das respostas deixadas em branco, com 7,3 \% dos estudantes, que se caracteriza como um máximo local. 6.3.2.4 Comentários sobre as respostas à Questão Discursiva 4 A Questão 4 foi sobre o aplicativo SOS Mulher e a sua utilidade para as mulheres que se encontram em situação de risco iminente.
A questão continha oito telas de interface desse aplicativo e solicitava -se uma avaliação de sua arquitetura da informação, uma exposição sobre o uso funcional das cores nas telas dessa interface e uma discussão sobre sua usabilidade.
Na questão, abordou -se uma temática contemporânea na área e tratou -se de uma das atividades profissionais das que possuem maior demanda no atual mercado de trabalho : o design digital.
No enunciado foram abordados pontos importantes das diretrizes curriculares — ergonomia informacional e desenvolvimento de interfaces para mídias digitais —, o que serviu para conduzir o estudante a uma reflexão sobre aspectos sensoriais e perceptivos que interferem na acessibilidade e usabilidade dos projetos de aplicativos para dispositivos móveis, correlacionando -as ao ato projetual e dando ênfase ao trabalho de organização dos elementos visuais de um projeto de design gráfico para meio digital.
O nível de conhecimento exigido para uma resposta que atendesse ao comando do enunciado é o necessário para a boa atuação de um profissional da área no mercado de trabalho mais amplo atualmente, uma vez que o Design Gráfico para meios digitais é hoje uma das áreas de maior demanda para a atuação dos profissionais do setor, com tendência de ampliação em função do crescimento do uso das novas tecnologias na sociedade brasileira, uma das mais ativas na internet.
Dentre as três questões discursivas propostas nesta edição do Exame, considerou -se essa a de maior grau de atualidade em relação ao exercício da profissão, uma vez que foi focada em aspectos extremamente contemporâneos das atividades no campo do Design Gráfico ; apesar disso, considerou -se o grau de dificuldade o mais baixo dentre todas as discursivas, pois a temática era muito presente na vida de estudantes de nível superior.
Não se pode deixar de mencionar que a parte mais trabalhosa para a elaboração das respostas foi o nível de detalhamento exigido no comando da questão.
O enunciado da questão funcionou de forma didática, com uma introdução que descrevia o objeto de análise, seguido de exemplo visual bem estruturado e detalhado, para só então colocar o comando específico, que foi formulado de uma maneira tal que orientava a ação do estudante para uma estrutura de resposta que facilitasse um desempenho satisfatório.
As imagens que compunham o enunciado possuíam papel central para a elaboração das respostas.
O enunciado da questão era claro e objetivo, com sequência de imagens com simulação da navegação de um aplicativo e, em seguida, comandos distintos para três itens : análise da arquitetura da informação, dissertação sobre o uso de cores e avaliação da usabilidade.
Para tal, no item ‘ a ’, considerou -se a análise do tipo de estrutura compositiva ( grid ), a avaliação da funcionalidade da composição centralizada e da distribuição assimétrica das áreas de acesso aos comandos, bem como da adequação da ocupação de tela e do dimensionamento dessas áreas para os diferentes itens e o design dos ícones em si e suas possibilidades de interpretação de seus significados em correlação com os comandos que designam.
No item ‘ b ’, foram avaliadas as considerações sobre a paleta de cores, sua adequação cultural ao público ao qual se destina o aplicativo, o contraste entre as cores utilizadas e o quanto facilitariam o uso das diferentes telas do aplicativo, bem como a correção do emprego de diferentes matizes ao melhor funcionamento do aplicativo.
No item ‘ c ’, foram consideradas as análises no que tange à usabilidade do aplicativo, tendo como guias fundamentalmente a jornada do usuário, a navegação, a legibilidade e a “ leiturabilidade ”, dentre outras. termos possuem significados muito próximos na área.
Nesse aspecto, a avaliação foi realizada com bastante parcimônia, levando -se em consideração em que pontos essas trocas estavam na avaliação e análise de pontos que são comuns às duas áreas ( Arquitetura de Informação e Usabilidade ) e quais realmente apresentaram uma confusão clara de conceito entre os dois itens.
Sobre o segundo item, a maior porcentagem de alunos avaliou a questão do uso das cores lilás ou magenta ( foram consideradas as possíveis alterações nas cores das impressões da prova, naturais em impressos desse tipo ).
Nesse item, foram consideradas tanto as respostas em que se levaram em consideração os elementos culturais quanto as respostas em que os estudantes questionaram de forma bem fundamentada a estereotipagem que o uso dessas cores pode gerar na sociedade.
Uma menor, mas expressiva porcentagem de estudantes, trouxe a questão do uso das cores contrastantes como forma de chamar a atenção.
Algumas respostas em que se apontaram críticas ao aplicativo ( com devido embasamento ) foram consideradas, por conterem outros pontos de vista, como por exemplo : o fato de o aplicativo usar a cor vermelha para destacar os comandos mais importantes, mas aparecer em momentos distintos, com funções positivas e negativas ( como em uma tela ela ser usada para “ mandar aviso para uma pessoa “ e em outra tela “ cancelar ” um comando ), o que poderia gerar confusão para o usuário.
Muito frequentemente foi possível encontrar respostas contendo apenas a descrição do aplicativo, com detalhes visuais de cada tela e com relatos óbvios das funções de cada botão, sem, contudo, um aprofundamento na análise da questão da arquitetura ou da usabilidade.
Outras vezes os alunos responderam somente sobre a importância social do aplicativo, fazendo uma análise da atual situação de violência contra a mulher na sociedade, ignorando a parte relativa ao design em si, não levando em consideração os aspectos pedidos nos itens ‘ a ’, ‘ b ’ e ‘ c ’ do enunciado.
Na imensa maioria das respostas a essas questões, não houve observância do que se solicitava no comando, e os estudantes se basearam apenas na análise de fundo sociológico ou psicológico sem menção a temas como Usabilidade ou Arquitetura da Informação.
Também foi possível encontrar respostas com discursos aplicados ao Design e a Usabilidade, mas sem abordagem do que estava exposto no enunciado e direcionado ao aplicativo em questão, que deveria servir de objeto de análise, de forma que não havia conexão entre o conhecimento do aluno e os comandos da questão.
Muitos erros comuns em relação ao uso das cores decorreram de previsões equivocadas do que poderia acontecer ao usuário.
Por exemplo : em vez de colocar quais as sensações que a cor roxa transmitia, de acordo com a Psicologia das Cores, o estudante afirmou que a cor foi bem aplicada porque deixava a mulher mais tranquila e relaxada quando em situação de perigo.
O que obviamente não caberia como análise estrita acerca de quais são as funções projetuais da utilização de tais cores.
Outro erro comum foram os textos com afirmativas de que não se pode avaliar concretamente o aplicativo, por estar realizando a atividade em imagens existentes em um impresso ( a prova ), quando o correto deveria ser exclusivamente o de avaliar no celular.
Outro ponto equivocado repetido com certa frequência foram as respostas cujas diretrizes apresentaram como referências sistemas operacionais e aplicativos no lugar de conceitos e autores.
Assim, o estilo “ Windows ” ou o “ Google Material ” foram muitas vezes citados como referências que deveriam ser seguidas.
No tocante às respostas avaliadas, pode -se notar que os cursos têm obtido sucesso na conscientização da importância de uma navegação de fácil compreensão para o usuário no desenvolvimento de sites e aplicativos, embora ainda se encontre carência em alguns pontos.
Poucas estudantes buscaram apontar exemplos de aplicativos como referência, e assim como poucos mencionaram aspectos teóricos, como as Heurísticas de Nielsen ou a Psicologia das Cores.
Aparentemente, são os itens que precisariam ser mais bem trabalhados nos cursos de Tecnólogos em Design Gráfico.
Outro ponto de destaque sobre a questão das cores foi a ausência de termos técnicos, que foram com frequência substituídos por uso de adjetivos, com o emprego de vocábulos que expressavam juízos de valor e adjetivações, e não ações projetuais objetivas.
Os estudantes empregaram em grande percentual expressões como “ cores calmas ”, “ layout feliz ”, “ design bonito e feminino ”, demonstrando certa distância entre o estudante e seu conhecimento dentro do vocabulário próprio de sua área de atuação.
Faz -se necessário, talvez, não só retrabalhar esses conceitos como conscientizar melhor o estudante sobre a importância do uso de termos técnicos para o exercício profissional, o que serve para fazer com que os clientes compreendam melhor os aspectos técnicos dos produtos de forma a não gerar confusão entre a atividade do design gráfico e a mera aplicação de gostos pessoais e/ou uso de elementos visuais como opções subjetivas.
Dentre o uso inadequado dos termos técnicos, destacou -se a importância da revisão do que é Arquitetura da Informação e do que é Usabilidade com a especificação de quando uma é consequência da outra, quais são os aspectos que caracterizam cada um desses termos e quais são os comuns aos dois, e com o entendimento de que a Arquitetura da Informação é a preocupação na disposição e diagramação das informações numa superfície ( seja ela digital ou impressa ) e que a Usabilidade é a relação que existe entre o usuário e a Além do pouco uso dos termos técnicos, foi possível identificar um grande problema quanto à qualidade de boa parte da redação das respostas.
Muitas apresentavam textos incompreensíveis ou com erros graves de ortografia e gramática.
É necessário reforçar trabalhos com leitura, interpretação de texto e redação dentro das universidades, sobretudo levando -se em conta que os futuros designers precisarão lidar com clientes, fazer apresentações, montar defesas de projeto.
Nada disso poderá ser desenvolvido sem o domínio da língua portuguesa. 6.3.2.5 Análise da Questão Discursiva 5 do Componente de Conhecimento Específico Na Tabela 6.20, constam as informações relativas à questão 5 do conjunto do Componente de Conhecimento Específico.
A nota Média dos estudantes de todo o Brasil foi 22,1.
A maior Média foi registrada na região Sul ( 26,6 ), enquanto a menor Média foi registrada na região Norte ( 16,6 ).
Quanto à variabilidade das notas, o Desvio padrão dos alunos do Brasil, como um todo, foi 21,1.
Enquanto o maior desvio foi encontrado na região Sul ( 23,4 ), o menor foi encontrado na região Norte ( 18,5 ).
A Mediana para o Brasil foi 20,0, a mesma obtida nas regiões Nordeste e Sudeste.
Nas regiões Norte e Centro-Oeste, a Mediana foi menor, 10,0 em ambas, enquanto na região Sul, a Mediana foi maior, 25,0.
Para o conjunto de alunos de Tecnologia em Design Gráfico do Brasil, a nota Máxima foi 100,0, obtida nas regiões Nordeste, Sudeste, Sul e Centro-Oeste.
A região Norte obteve nota Máxima 70,0.
A nota Mínima do Brasil foi zero, o mesmo valor encontrado em todas as regiões.
No Gráfico 6.13, é representada a distribuição das notas da questão discursiva 5 do Componente de Conhecimento Específico.
Essa distribuição tem moda no intervalo correspondente aos que tiraram nota zero, com 19,7 \% do total.
Somando -se a frequência desse intervalo com a daqueles que deixaram a resposta em branco chega -se próximo a 30 \% dos participantes. 6.3.2.6 Comentários sobre as respostas à Questão Discursiva 5 O conteúdo da questão tinha como foco uma temática clássica na área de Design Gráfico ( Sistemas de Identidade Visual ) e nela trabalhava -se conhecimentos que agregavam aspectos de Gerenciamento e de Avaliação de Viabilidade Técnica, funcional e econômica desse tipo de produto gráfico na sociedade e no mercado de atuação dos clientes.
No enunciado, abordavam -se esses importantes pontos das diretrizes curriculares de forma bem estruturada, com embasamento em referência bibliográfica consagrada na área e com exemplificação visual pertinente.
No comando da questão solicitava -se que o aluno discorresse sobre o redesign de marcas de empresas tradicionais, abordando a preservação da identidade e a fidelização da marca.
O nível de conhecimento exigido para uma boa resposta é o necessário para a boa atuação profissional no mercado de design gráfico, em especial no desenvolvimento de tarefa das mais constantes e tradicionais, uma vez que identidades visuais são trabalhos típicos e cotidianos para os profissionais do setor, sendo certamente uma das atividades de maior incidência na vida laboral dos designers gráficos.
O enunciado da questão estava constituído a partir de uma reflexão teórica, seguida de exemplificação visual, após a qual o comando foi formulado de maneira concisa, objetiva e estruturada, não deixando margem para interpretações divergentes.
Essa questão foi a que abordou a temática mais tradicional do exercício da profissão, dentre as três que compunham a prova, complementando de forma bastante equilibrada o espectro de conteúdos sobre as atividades do campo do Design Gráfico.
O grau de dificuldade da questão foi elevado, pois exigiu que conhecimentos basilares, de mais de uma atividade, fossem articulados para o atendimento dos três tópicos presentes no comando do enunciado.
A questão foi formulada em linguagem direta e objetiva, e o enunciado abordou conteúdo relevante para a formação profissional do estudante de forma clara, aplicando -o a um exemplo da prática profissional dos mais comuns na rotina de designers gráficos, uma vez que a análise demandada é de caráter geral sobre o tema de retrabalho e gestão de identidades visuais.
Houve concisão e precisão nos comandos e facilidade para a apreensão dos aspectos que a resposta deveria contemplar, com margem muito pequena de variação das interpretações possíveis do enunciado, inclusive com o exemplo exposto direcionando a resposta para os casos de empresas que necessitam manter vínculos fortes entre as identidades visuais antigas e as que resultam do processo de redesign.
Os comandos foram claros e solicitava -se que o estudante dissertasse sobre o processo de redesign em marcas tradicionais, e não genericamente, pedia -se que abordasse os aspectos visuais que permitissem a criação de vínculos entre as visualidades envolvidas no processo para a preservação da identidade, através dos quais há o reconhecimento da marca, o que propiciaria a manutenção da fidelização dos consumidores, uma vez que deveria ser gerada a percepção de que a nova identidade visual é uma atualização da marca original do produto ou serviço.
Portanto, era esperado que o estudante : descrevesse o processo de redesign ( ainda que sinteticamente ) e avaliasse seu emprego em marcas tradicionais como forma de atualizá -las sem perda dos vínculos identitários ; elencasse os elementos que compunham uma identidade visual e analisasse como esses deveriam ser visualmente trabalhados no processo para que fosse bem sucedido ; e, por fim, discorresse sobre a relação dessas soluções visuais com a parcela do público já fidelizado pela marca, podendo ainda se estender tratando dos impactos da ação sobre novos públicos.
A primeira parte ( descrição do processo de redesign ) foi a com maior frequência de acerto, ainda que de forma pouco estruturada e sem a análise demandada.
Nesse aspecto, foi avaliado que o próprio enunciado provocava nos estudantes um estímulo no sentido de recordar o conteúdo ministrado na graduação e os levava a desenvolver o raciocínio de forma correta, ainda que incompleta.
Os elementos constitutivos de uma identidade visual e o seu tratamento no processo de redesign foram os aspectos que apresentaram o menor índice de acerto, sendo frequentemente omitidos nas respostas e com constantes erros na utilização do vocabulário técnico consagrado, ainda que os considerando de forma condescendente com relação ao uso mais rigoroso.
Foram frequentes as menções a mascotes, slogans, anúncios, publicidade e outros itens equivocados como sendo elementos componentes de identidades visuais.
No tocante às mascotes, quando citadas em associação aos exemplos em que poderiam ser consideradas como símbolos figurativos e componentes de identidades corporativas ou de eventos, a citação foi pontuada ; mas nos casos de menções soltas, sem exemplificação que as contextualizasse, a pontuação não foi concedida.
Respostas em que fossem contempladas as estratégias e recomendações para o bom redesign dos elementos das identidades visuais foram os pontos menos desenvolvidos pelos estudantes, sendo preocupante o elevado percentual daqueles que nem ao mesmo citaram uma.
Dentre os elementos compositivos de identidades visuais, mesmo o mais óbvio, o logotipo, foi raramente citado.
Também foi frequente o uso da palavra “ marca ” para se referir às assinaturas visuais, o que foi considerado correto nos casos em que os estudantes apresentaram um desenvolvimento acertado e apenas empregaram o termo com os sentidos de identidade visual, assinatura visual ou logotipo.
Cabe ressaltar que em parte significativa das notas máximas, os estudantes utilizaram como estratégia a inclusão da descrição das imagens presentes no enunciado e, a partir dessa descrição, desenvolveram argumentações que efetivamente atendiam aos comandos do enunciado.
Outra estratégia frequente que resultou em baixa pontuação foi a descrição sintética do processo de redesign e a citação dos vocábulos referentes aos elementos constitutivos de Identidades Visuais ( logotipo, símbolo ou cores institucionais ), mas sem a análise do primeiro ou considerações sobre o trabalho a ser realizado com os segundos.
Uma grande parte dos alunos que conseguiu descrever a análise do processo e os aspectos da influência de um bom redesign na manutenção da fidelização da marca, mas sem fazer referência aos elementos gráficos das identidades visuais, obteve pontuações medianas.
Menos frequente foi o quantitativo de alunos que associou a resposta correta do primeiro tópico ( redesign ) a um desempenho bom no que diz respeito aos elementos visuais constitutivos de uma identidade visual.
Igualmente pequeno foi o quantitativo de estudantes alcançou pontuação mediana por apresentar respostas parciais aos três tópicos.
Muito raros foram os estudantes que apresentaram um bom desenvolvimento sobre os elementos das identidades visuais e, por conseguinte, obtiveram pontuação parcial em um dos outros dois tópicos.
As estratégias que levaram às melhores avaliações foram diversificadas, desde as abordagens muito calcadas em exemplificações ( com destaque para o caso exposto nas imagens do enunciado ) até respostas em que os estudantes que contemplaram os comandos de forma absolutamente abstrata, trabalhando com formulação teórica bem consolidada e conceituação corretamente aplicada a todos os tópicos exigidos.
Como o maior percentual de sucesso foi pertinente ao conteúdo mais afeito à interface Design/Marketing, tanto no tocante ao redesign, quanto ao tópico de Fidelização, a conclusão mais coerente, em face das respostas dos alunos, foi de que esse aspecto foi mais bem trabalhado e/ou absorvido pelos estudantes.
Já o aspecto mais próximo ao campo do design em si ( elementos da Identidade Visual ) foi exatamente o que apresentou menor índice de acerto nas respostas.
Tal fato apontou para um baixo grau de absorção dos conteúdos mais basilares do campo do Design Gráfico, em que pese sua presença consistente nas diretrizes e nas grades curriculares, além de uma notável dificuldade com o vocabulário profissional mais cotidiano e com o desenvolvimento de respostas que articulem com desenvoltura a associação entre teoria e prática, que o enunciado da questão claramente visava a avaliar.
Na avaliação, foi considerado se o estudante descreveu o processo de redesign como a atualização estilística dos elementos visuais de uma marca, sem que houvesse perda de sua identidade e a consequente desvinculação mnemônica dos principais elementos de sua visualidade, o que poderia interferir na fidelização de marca já estabelecida.
Para tal, foi necessário manter as referências de aspectos como a cor institucional, as características gerais do design do logotipo, bem como as da tipografia/tipologia empregada, e da visualidade do símbolo.
Esses elementos foram avaliados e os aspectos descritos foram pensados em função de se conseguir que a marca tivesse uma imagem mais atualizada com o mercado, porém sem perder seu potencial de reconhecimento para não interferir na fidelização de marca, pois alterações muito grandes poderiam gerar uma identidade visual completamente desvinculada do trabalho original, o que evocaria um novo produto ou serviço, criando, com isso, dificuldades para que o público já fidelizado identificasse o produto ou serviço através desses elementos.
Tendo como base o descrito no parágrafo anterior, com a correção, ficou evidente a falta de familiaridade de grande percentual dos estudantes com os tópicos.
O destaque negativo ficou por conta exatamente do ponto central para o Design Gráfico : a materialidade das identidades visuais, seus elementos gráficos, sua nomenclatura e as características que regem o trabalho para a obtenção de um design com elevado grau de efetividade no mercado, este último sendo o ponto com mais elevado grau de omissão.
O ponto que mais chamou atenção foi o baixo aproveitamento no aspecto mais fundamental da questão com relação ao Design Gráfico : pouca articulação entre a questão das identidades visuais e seus elementos constitutivos.
O baixo rendimento apresentado no uso do jargão profissional foi verdadeiramente preocupante, visto que termos como logotipo, símbolo, cores institucionais, assinaturas visuais, dentre outros são sabidamente trabalhados em todos os cursos.
Na solução dessa questão, ficou evidente a não absorção desse vocabulário pela imensa maioria dos estudantes e a não vinculação dele ao trabalho de redesign de marcas.
As respostas dos estudantes evidenciaram um alarmante problema com o uso da língua portuguesa e com a redação de respostas organizadas, coesas e coerentes.
O percentual de respostas desconexas e/ou com graves problemas de estruturação lógica foi bastante significativo, somado às respostas em que havia dificuldade de articulação da expressão escrita e às redigidas de forma tópica e incompleta a apenas um dos três itens solicitados.
Em face à correção da questão discursiva 5, seria razoável que as coordenações dos cursos de tecnólogo em Design Gráfico intensificassem exercícios para a melhoria da capacidade de redação dos estudantes e também do domínio do jargão profissional, uma vez que a vida profissional exigirá a apresentação de material escrito a clientes em várias oportunidades. 6.3.3 Considerações Finais Em relação à questão discursiva 3, não houve diversidade nem originalidade nas respostas, além de nem todas terem apresentado propostas de projetos ou terem feito conexões com conceitos de História da Arte.
Boa parte dos alunos usou vocabulário pouco correlacionado ao tema, sem, ao menos, conseguir fazer conexões concretas com projetos gráficos existentes ou hipotéticos, simplificando as possibilidades de resposta ao conjunto de dados mais presentes na cultura de massa, aos estereótipos da africanidade subsaariana ou a dados das religiões de matriz africana.
Uma parte significativa dos alunos parafraseou o texto ou simplesmente apresentou como solução itens presentes no próprio texto.
As respostas dos alunos demonstraram o profundo desconhecimento de elementos básicos de projetos gráficos, do continente africano e de suas culturas.
Foram identificadas muitas deficiências em relação ao desenvolvimento das argumentações para a construção de respostas adequadas ao comando da Questão 3, aliadas a um baixo domínio da teoria da Área de Tecnologia em Design de Interiores.
Um fato alarmante foi a constatação de que, em um curso com objetivo de formar profissionais capazes de pensar projetos gráficos novos e diversificados, há grande dificuldade na compreensão da exigência de relação da aplicação dos elementos em um projeto gráfico, concretamente existente ou hipotético.
A Questão 4 versou sobre um tema mais atual da Área de Tecnologia em Design Gráfico, por conta da abordagem do conhecimento necessário para o bom desempenho desse profissional no mercado de trabalho, mercado esse cada vez mais amplo e com uma grande demanda de profissionais qualificados e atualizados com o uso de novas tecnologias.
Mas não se pode deixar de destacar que em poucas respostas havia exemplos de aplicativos como referência.
Escassos eram os apontamentos sobre aspectos teóricos, como as Heurísticas de Nielsen ou a Psicologia das Cores, o uso de termos técnicos.
Também foi possível identificar um grande problema quanto à qualidade de boa parte da redação das respostas.
A questão discursiva 5 abordou uma temática mais tradicional do exercício da profissão, quando comparada às outras duas que compunham a prova, complementando, de forma bastante equilibrada, o espectro de conteúdos sobre as atividades do campo do Design Gráfico abrangido pelas questões discursivas, uma vez que a análise demandada foi de caráter geral sobre o tema de Retrabalho e Gestão de Identidades Visuais.
Na maioria das respostas, a conexão entre a manutenção de aspectos visuais característicos da marca pela nova identidade visual e o vínculo mnemônico com a anterior foi pouco citado, ou citado de forma muito sucinta.
Foram encontradas, também, respostas em que os estudantes associaram a descrição sintética do processo de redesign ao fato de que alguma manutenção de vínculo visual com a identidade visual antiga deveria ser determinante na identificação do produto ou serviço pelo cliente já fidelizado.
Pelas respostas dos alunos, verificou -se um baixo grau de absorção dos conteúdos mais basilares do campo do Design Gráfico, uma grande dificuldade com o vocabulário profissional mais cotidiano, com o desenvolvimento de respostas que articulem com desenvoltura a associação entre teoria e prática e, também, pouca articulação entre a questão das identidades visuais e seus elementos constitutivos.

% matched lemmas: laranja, traçar
\end{textsample}
